\exercisegroup

\begin{enumerate}
\def\labelenumi{\arabic{enumi})}
\item
  Let X be a Lindelöf space (TG, IX, appendix 1, p. 75, def. 1) that is loopable. Show that, for every point \(x \in X\) , the Poincaré group \(\pi_1(X, x)\) is countable.
\item
  Let Y be a uniform space and let X be a subspace of Y. Let \((\mathrm{U}_{i})_{i\in\mathrm{I}}\) be an open cover of X.
\end{enumerate}

\begin{enumerate}
\def\labelenumi{\alph{enumi})}
\tightlist
\item
  Prove that there exists an open neighborhood V of X in Y and an open covering \((V_{i})_{i\in I}\) of V satisfying the following properties:
\end{enumerate}

\begin{enumerate}
\def\labelenumi{(\roman{enumi})}
\item
  For every \(i\) , \(\mathrm{U}_i = \mathrm{V}_i \cap \mathrm{X}\) ;
\item
  For every subset J of I such that \(\bigcap_{j\in J}V_j\neq \varnothing\), we have \(\bigcap_{j\in J}U_j\neq \varnothing .\)
\end{enumerate}

\begin{enumerate}
\def\labelenumi{\alph{enumi})}
\setcounter{enumi}{1}
\item
  Assume that Y is locally connected and that the \(U_{i}\) are connected. Prove that there exists such a covering in which the \(V_{i}\) are connected open subsets of Y.
\item
  Suppose that the \(U_{i}\) are connected and let, for every \(i \in I\), a point \(x_{i} \in U_{i}\). Suppose that the image of \(\pi_{1}(U_{i}, x_{i})\) in \(\pi_{1}(X, x_{i})\) is reduced to the identity element. Show that there exists a neighborhood V of X in Y such that, for every \(x \in X\), the canonical homomorphism \(\pi_{1}(X, x) \to \pi_{1}(V, x)\) admits a retraction.
\end{enumerate}

\begin{enumerate}
\def\labelenumi{\arabic{enumi})}
\setcounter{enumi}{2}
\tightlist
\item
  \begin{enumerate}
  \def\labelenumii{(\arabic{enumii})}
  \setcounter{enumii}{2}
  \tightlist
  \item
    Let \(X\) be a topological space and let \(x\) be a point of \(X\) .
  \end{enumerate}
\end{enumerate}

Let G be a group and let \(\varphi\colon\pi_{1}(\mathrm{X},x)\to\mathrm{G}\) be a group homomorphism. Let \((\mathrm{U}_{n})_{n\in\mathbf{N}}\) be a decreasing sequence of subspaces of X whose intersection is reduced to the point \(x\) . For all \(n\) , let \(\mathrm{G}_{n}\) be the set of elements of G that are images under \(\varphi\) of the class of a loop contained in \(\mathrm{U}_{n}\) .

\begin{enumerate}
\def\labelenumi{\alph{enumi})}
\item
  If G is countable, the sequence \((\mathrm{G}_{n})\) is stationary. (Reduce to the case where \((\mathrm{G}_{n}:\mathrm{G}_{n+1})\geqslant n+1\) for all n and consider an enumeration \((\gamma_{n})_{n}\) of G; construct recursively elements \(\lambda_{n}\in G_{n}\) such that, on writing \(w_{n}=\lambda_{1}\ldots\lambda_{n}\), one has \(\gamma_{i}\notin w_{n}G_{n+1}\) for \(1\leqslant i\leqslant n\). Show then that the intersection \(\bigcap_{i}w_{i}G_{i+1}\) is not empty.)
\item
  Suppose that G is abelian and that its identity element is the only infinitely divisible element of G. Then \(\bigcap_{n} G_{n}\) is reduced to the identity element. (Let a be an element of the intersection; for every n, let \(\gamma_{n} \in \pi_{1}(U_{n}, x)\) whose image under \(\varphi\) is equal to a. Construct by induction path classes \(\lambda_{n} \in \pi_{1}(U_{n}, x)\) such that \(\lambda_{n} = \gamma_{n}\lambda_{n+1}^{d}\), where d is a fixed integer \(\geqslant 2\). Show that \(b = (d - 1)\varphi(\lambda_{1}) + a\) is infinitely divisible. Deduce that a = 0.)
\item
  More generally, if G is abelian and every infinitely divisible element of G is a torsion element, then the intersection of the \(G_{n}\) consists of torsion elements.
\end{enumerate}

(3) The results of this exercise and the following one are taken from the article by J. W. CANNON and G. R. CONNER, “On the fundamental groups of one-dimensional spaces,” Topology and its Applications 153 (2006), pp. 2648–2672.

\begin{enumerate}
\def\labelenumi{\arabic{enumi})}
\setcounter{enumi}{3}
\tightlist
\item
  Let X be a connected and locally connected by arcs topological space. Let x be a point of X that has a countable fundamental system of neighborhoods. Let S be a set and let \(\varphi\colon\pi_{1}(\mathrm{X},x)\to\mathbf{Z}^{(\mathrm{S})}\) be a surjective group homomorphism.
\end{enumerate}

\begin{enumerate}
\def\labelenumi{\alph{enumi})}
\item
  Prove that S is countable.
\item
  Assume that X is compact; prove that S is finite.
\item
  We assume that \(\pi_{1}(X,x)\) is a free abelian group; prove that X is loopable.
\item
  We assume that \(\pi_{1}(X,x)\) is a free group; prove that X is loopable. (Use the fact that the intersection of the subgroups of the lower central series of a free group is reduced to the identity element.)
\end{enumerate}

\begin{enumerate}
\def\labelenumi{\arabic{enumi})}
\setcounter{enumi}{4}
\tightlist
\item
  Let P be the Hawaiian earring (III, p. 336, exerc. 6, whose notation we adopt). For \(n \in \mathbf{N}^*\) , let \(\alpha_n \in \pi_1(\mathrm{P}, o)\) be the class of a loop in \(\mathrm{C}_n\) whose class generates \(\pi_1(\mathrm{C}_n, 0)\) .
\end{enumerate}

\begin{enumerate}
\def\labelenumi{\alph{enumi})}
\item
  Prove that the map \(\mathrm{Hom}(\pi_{1}(\mathrm{P},0),\mathbf{Z})\to\mathbf{Z}^{\mathbf{N}}\) , \(\varphi\mapsto(\varphi(\alpha_{n}))_{n\in\mathbf{N}}\) is an injective group homomorphism and that its image is equal to \(\mathbf{Z}^{(\mathbf{N})}\) . (4)
\item
  The group \(\pi_{1}(P,0)\) is not a free group; the quotient of \(\pi_{1}(P,0)\) by its derived group is not a free abelian group.
\item
  Let A be the Hawaiian archipelago (III, p. 338, exerc. 10). Prove that every homomorphism from \(\pi_1(A, o)\) into \(\mathbf{Z}\) is trivial.
\end{enumerate}

(4) See B. DE SMIT, “The fundamental group of the Hawaiian earring is not free,” International Journal of Algebra and Computation, Vol. 2, No. 1 (1992), pp. 33–37.

\exercisegroup

\begin{enumerate}
\def\labelenumi{\arabic{enumi})}
\item
  Let G be the topological group O(2, R). We have \(\pi_{0}(G) \simeq Z/2Z\) and \(\pi_{1}(G) \simeq Z\) ; the nontrivial element of \(\pi_{0}(G)\) acts on \(\pi_{1}(G)\) by the mapping \(t \mapsto -t\) .
\item
  We denote by H the algebra of Hamilton quaternions (TG, VIII, p. 4) and \((1, i, j, k)\) its canonical basis. Let N denote the norm quadratic form of H; the set of quaternions of norm 1 is the sphere \(S_{3}\) . Let \(H_{0}\) be the real vector subspace of quaternions with zero trace and \(N_{0}\) the restriction to \(H_{0}\) of N.
\end{enumerate}

\begin{enumerate}
\def\labelenumi{\alph{enumi})}
\item
  Prove that the homomorphism in Proposition 4 of TG, VIII, p. 4 makes the sphere \(\mathbf{S}_3\) a covering of degree 2 of \(\mathbf{SO}(\mathrm{N}_0)\) .
\item
  Hence deduce that the Poincaré group of SO(3, R) is isomorphic to Z/2Z. Explicitly describe a loop in \(I_{3}\) in SO(3, R) whose strict homotopy class is the unique nontrivial element of this Poincaré group.
\item
  Let \(\varphi\) be the map from \(\mathbf{S}_3 \times \mathbf{S}_3\) into SO(4, R) given by \(\varphi(\mathbf{q}_1, \mathbf{q}_2)(\mathbf{x}) = \mathbf{q}_1 \mathbf{x} \mathbf{q}_2^{-1}\). Prove that \(\varphi\) makes \(\mathbf{S}_3 \times \mathbf{S}_3\) a covering of degree 2.
\item
  Deduce that the Poincaré group of SO(4, R) is isomorphic to Z/2Z.
\end{enumerate}

\begin{enumerate}
\def\labelenumi{\arabic{enumi})}
\setcounter{enumi}{2}
\tightlist
\item
  Let n be a natural number. Let the topological group \(\mathrm{SO}(n+1,\mathbf{R})\) act on the sphere \(S_{n}\) of \(R^{n+1}\) .
\end{enumerate}

\begin{enumerate}
\def\labelenumi{\alph{enumi})}
\item
  Prove that the fixator of the vector \(e_{n+1}\) is identified with the topological group \(\mathrm{SO}(n,\mathbf{R})\) .
\item
  Let \(f \colon \mathbf{I} \to \mathbf{S}_n\) be a continuous map. Assume that \(f(t) \neq (0, \ldots, 0, 1)\) for all \(t \in \mathbf{I} .\) Prove that there exists a continuous map \(g \colon \mathbf{I} \to \mathrm{SO}(n + 1, \mathbf{R})\) such that \(f(t) = g(t) \cdot \mathbf{e}_{n+1}\) for all \(t \in \mathbf{I} \).
\item
  If \(n \geqslant 3\) , prove that the inclusion of \(\mathrm{SO}(n, \mathbf{R})\) into \(\mathrm{SO}(n+1, \mathbf{R})\) induces an isomorphism between Poincaré groups.
\end{enumerate}

\begin{enumerate}
\def\labelenumi{\arabic{enumi})}
\setcounter{enumi}{3}
\tightlist
\item
  Let n be a natural number. Let \(B_{0}\) be the subgroup of \(\mathbf{SL}(n,\mathbf{R})\) consisting of upper triangular matrices whose diagonal coefficients are all strictly positive.
\end{enumerate}

\begin{enumerate}
\def\labelenumi{\alph{enumi})}
\item
  Prove that the map from \(\mathrm{SO}(n,\mathbf{R})\times\mathrm{B}_{0}\) into \(\mathbf{SL}(n,\mathbf{R})\) given by \((g,u)\mapsto g\cdot u\) is a homeomorphism (cf. INT, VII, p. 91, prop. 7).
\item
  Prove that the inclusion of SO(n, R) into SL(n, R) induces an isomorphism between Poincaré groups.
\end{enumerate}

\begin{enumerate}
\def\labelenumi{\arabic{enumi})}
\setcounter{enumi}{4}
\tightlist
\item
  Let G be a connected loopable topological group, let \((\widetilde{\mathrm{G}}, p)\) be a universal cover of G; we endow \(\widetilde{G}\) with the topological group structure for which p is a group homomorphism. We denote by Z the center of G, by e its identity element, \(\widetilde{Z}\) the center of \(\widetilde{G}\) and \(\widetilde{e}\) its identity element.
\end{enumerate}

\begin{enumerate}
\def\labelenumi{\alph{enumi})}
\item
  Let \(f\) be a continuous automorphism of the group G. Prove that there exists a unique continuous map \(\widetilde{f} \colon \widetilde{\mathrm{G}} \to \widetilde{\mathrm{G}}\) such that \(\widetilde{f}(\widetilde{e}) = \widetilde{e}\) and \(p \circ \widetilde{f} = f \circ p\). Prove that \(\widetilde{f}\) is a group automorphism. Prove that the map \(\varphi \colon f \mapsto \widetilde{f}\) is an injective group homomorphism from \(\operatorname{Aut}(\mathbf{G})\) into \(\operatorname{Aut}(\widetilde{\mathbf{G}})\) .
\item
  We denote \(\operatorname{Out}(\mathbf{G})\) the quotient of the group \(\operatorname{Aut}(\mathbf{G})\) by the subgroup of inner automorphisms; one defines \(\operatorname{Out}(\widetilde{\mathbf{G}})\) analogously. Prove that the homomorphism \(\varphi\), induced by passing to quotients, is an injective group homomorphism \(\psi\colon\operatorname{Out}(\mathbf{G})\to\operatorname{Out}(\widetilde{\mathbf{G}})\) .
\end{enumerate}

\begin{enumerate}
\def\labelenumi{\arabic{enumi})}
\setcounter{enumi}{5}
\tightlist
\item
  Let G and H be connected real Lie groups, and let \(f: G \to H\) be a surjective morphism of Lie groups. We denote by \(Z_{G}\) the center of G and by \(Z_{H}\) that of H.
\end{enumerate}

\begin{enumerate}
\def\labelenumi{\alph{enumi})}
\item
  Prove the equivalence of the following assertions: (i) The map f makes G a covering of H; (ii) \(\operatorname{Ker}(f)\) is discrete; (iii) The morphism \(\mathrm{L}(f)\colon\mathrm{L}(\mathrm{G})\to\mathrm{L}(\mathrm{H})\) induced from f by passing to Lie algebras is an isomorphism.
\item
  Prove that then \(f(\mathrm{Z}_{\mathrm{G}})=\mathrm{Z}_{\mathrm{H}}\) and that \(f^{-1}(\mathrm{Z}_{\mathrm{H}})=\mathrm{Z}_{\mathrm{G}}\) .
\end{enumerate}

\begin{enumerate}
\def\labelenumi{\arabic{enumi})}
\setcounter{enumi}{6}
\tightlist
\item
  Let G be a real Lie group, denote by \(G_{0}\) its neutral component, \(Z_{0}\) the center of \(G_{0}\) and \(\pi_{0}(G)=G/G_{0}\) . Let \((\widetilde{G}_{0},p_{0})\) be a universal covering of \(G_{0}\) , \(\widetilde{Z}_{0}\) the center of \(\widetilde{G}_{0}\) and \(\widetilde{e}\) its identity element.
\end{enumerate}

\begin{enumerate}
\def\labelenumi{\alph{enumi})}
\item
  Let \(\theta\colon\pi_{0}(\mathrm{G})\to\mathrm{Out}(\mathrm{G}_{0})\) be the homomorphism associated with the extension \(\mathcal{E}\colon\mathrm{G}_{0}\to\mathrm{G}\to\pi_{0}(\mathrm{G})\) (A, X, §7, exerc. 5). Prove that the cohomology class \(\omega(\pi_{0}(\mathrm{G}),\mathrm{Z}_{0},\theta)\) of \(\mathrm{H}^{3}(\pi_{0}(\mathrm{G}),\mathrm{Z}_{0})\) defined in loc. cit. is zero.
\item
  Assume that there exists a Lie group \(\widetilde{\mathrm{G}}\), a surjective map \(p\colon\widetilde{\mathrm{G}}\to\mathrm{G}\), and an isomorphism \(j\) of \(\widetilde{\mathrm{G}}_{0}\) onto the neutral component of \(\widetilde{\mathrm{G}}\) such that \(p_{0}=p\circ j\). Show that we have \(\omega(\pi_{0}(\mathrm{G}),\widetilde{\mathrm{Z}}_{0},\psi\circ\theta)=0\) in \(\mathrm{H}^{3}(\pi_{0}(\mathrm{G}),\widetilde{\mathrm{Z}}_{0})\) , where \(\psi\colon\mathrm{Out}(\mathrm{G}_{0})\to\mathrm{Out}(\widetilde{\mathrm{G}}_{0})\) is the homomorphism defined in exerc. 5.
\item
  Assume that \(\omega(\pi_{0}(\mathrm{G}),\widetilde{\mathrm{Z}}_{0},\psi\circ\theta)=0\) . Let \(\widetilde{E}\colon\widetilde{G}_{0}\to E\to\pi_{0}(G)\) be an extension of \(\pi_{0}(G)\) by \(\widetilde{G}_{0}\) whose associated homomorphism is equal to \(\psi\circ\theta\). Prove that \(\mathrm{E}/\pi_{1}(\mathrm{G}_{0})\) is an extension of \(\pi_{0}(G)\) by \(G_{0}\) whose associated homomorphism is equal to \(\theta\) .
\end{enumerate}

Now let \(\gamma\) be the unique element of \(\mathrm{H}^2 (\pi_0(\mathrm{G}),\mathrm{G}_0)\) such that \(\gamma \cdot [\mathrm{E} / \pi_1(\mathrm{G}_0)] = [\mathrm{G}]\) (loc. cit., e)). Consider the homomorphism \(\delta \colon \mathrm{H}^2 (\pi_0(\mathrm{G}),\mathrm{Z}_0)\to\) \(\mathrm{H}^3 (\pi_0(\mathrm{G}),\pi_1(\mathrm{G}))\) associated with the exact sequence of abelian groups \(1\to \pi_1(\mathrm{G}_0)\to\) \(\widetilde{\mathbf{Z}}_0\to \mathbf{Z}_0\to 1\) (exerc. 6). Show that the element \(\delta (\gamma)\) of \(\mathrm{H}^3 (\pi_0(\mathrm{G}),\pi_1(\mathrm{G}))\) does not depend on the choice of the extension \(\widetilde{\mathcal{E}}\) ; it is denoted \(\delta (\mathrm{G})\) .

d) Assume that \(\omega (\pi_0(\mathrm{G}),\widetilde{\mathrm{Z}}_0,\psi \circ \theta) = 0\) . For there to exist a Lie group \(\widetilde{\mathrm{G}}\), a surjective map \(p\colon \widetilde{\mathrm{G}}\to \mathrm{G}\), and an isomorphism \(j\) of \(\widetilde{\mathrm{G}}_0\) onto the neutral component of \(\widetilde{G}\) such that \(p_{0}=p\circ j\) , it is necessary and sufficient that \(\delta(\mathrm{G})=0.\) (5)

(5) For examples, see the article by R. L. TAYLOR, “Covering groups of nonconnected topological groups,” Proc. Amer. Math. Soc. 5 (1954), pp. 753–768.

\begin{enumerate}
\def\labelenumi{\arabic{enumi})}
\setcounter{enumi}{7}
\tightlist
\item
  Let G be a Lie group and let H be a closed subgroup of G; we denote by \(H_{0}\) the neutral component of H and \(\pi_{0}(H)\) the group \(H/H_{0}\) . Let us also denote by p the canonical surjection from G to G/H.
\end{enumerate}

\begin{enumerate}
\def\labelenumi{\alph{enumi})}
\item
  The map p has the path lifting property.
\item
  If H is connected, the homomorphism \(p_{*}\colon\pi_{1}(\mathrm{G},e)\to\pi_{1}(\mathrm{G}/\mathrm{H},p(e))\) is surjective.
\item
  There exists a unique map \(\varphi\colon\pi_{1}(\mathrm{G/H},p(e))\to\pi_{0}(\mathrm{H})\) which, for every path \(\widetilde{c}\) with source \(e\) in G whose term belongs to H, associates to the class of the loop \(p\circ\widetilde{c}\) in G/H the class of \(\widetilde{c}(1)\) modulo \(\mathrm{H}_{0}\) ; it is a group homomorphism.
\item
  If G is simply connected, then \(\varphi\) is an isomorphism. In particular, H is connected if and only if G/H is simply connected.
\end{enumerate}

\begin{enumerate}
\def\labelenumi{\arabic{enumi})}
\setcounter{enumi}{8}
\tightlist
\item
  Let G be a connected Lie group, let \((\widetilde{\mathbf{G}}, q)\) be a universal cover of G and let \(\widetilde{e}\) be the identity element of \(\widetilde{G}\). Let H be a closed connected subgroup of G; we denote by \(i: H \to G\) the canonical injection and \(p: G \to G/H\) the canonical surjection.
\end{enumerate}

\begin{enumerate}
\def\labelenumi{\alph{enumi})}
\item
  Let \(H_{1}\) be the neutral component of \(\overline{q}^{-1}(H)\) and let \(q_{1}\colon\widetilde{\mathrm{G}}/\mathrm{H}_{1}\to\mathrm{G}/\mathrm{H}\) be the map induced by q by passing to the quotients. Prove that \(q_{1}\) makes \(\widetilde{G}/H_{1}\) a covering of G/H whose fiber at \(p(e)\) is isomorphic to \(\pi_{1}(\mathrm{G}/\mathrm{H},p(e))\) .
\item
  Deduce that there exists an exact sequence of groups
\end{enumerate}

\[
\pi_ {1} (\mathrm{H} _ {1}, \widetilde {e}) \xrightarrow {q _ {*}} \pi_ {1} (\mathrm{H}, e) \xrightarrow {i _ {*}} \pi_ {1} (\mathrm{G}, e) \xrightarrow {p _ {*}} \pi_ {1} (\mathrm{G} / \mathrm{H}, p (e)).
\]

\begin{enumerate}
\def\labelenumi{\arabic{enumi})}
\setcounter{enumi}{9}
\tightlist
\item
  \begin{enumerate}
  \def\labelenumii{\alph{enumii})}
  \tightlist
  \item
    Prove that, for every integer \(n \geqslant 2\) , the group \(\mathrm{SU}(n, \mathbf{C})\) is simply connected. (The space \(\mathrm{SU}(2, \mathbf{C})\) is homeomorphic to \(S_{3}\). Let \(\mathrm{SU}(n, \mathbf{C})\) act on the unit sphere of \(C^{n}\), identify the stabilizer of the point \((0, \ldots, 0, 1)\) with the group \(\mathrm{SU}(n - 1, \mathbf{C})\), and use Exercise 9.)
  \end{enumerate}
\end{enumerate}

\begin{enumerate}
\def\labelenumi{\alph{enumi})}
\setcounter{enumi}{1}
\tightlist
\item
  Prove that for every integer \(n \geqslant 2\) , the group \(\mathbf{SL}(n, \mathbf{C})\) is simply connected. (Let \(B_{+}\) be the subgroup of \(\mathbf{SL}(n, \mathbf{C})\) consisting of upper triangular matrices whose diagonal coefficients are real numbers \textgreater{} 0. Show that the map from \(\mathrm{SU}(n, \mathbf{C}) \times \mathrm{B}_{+}\) into \(\mathbf{SL}(n, \mathbf{C})\) given by \((g, u) \mapsto g \cdot u\) is a homeomorphism.)
\end{enumerate}

\exercisegroup

\begin{enumerate}
\def\labelenumi{\arabic{enumi})}
\item
  Let X and Y be topological spaces, and let \(f: X \to Y\) be a continuous map. Let \((T, p)\) be a Y-space. The canonical descent datum \(\tau\) on the X-space \(Z = X \times_{Y} T\) is effective (IV, p. 385, prop. 1), the canonical map \(Z/R_{\tau} \to T\) is injective and continuous. It is not, however, necessarily surjective, nor strict.
\item
  Let Y be a topological space obtained by gluing spaces \(X_{i}\). Let X be the topological space that is the disjoint union of the family \((\mathrm{X}_{i})\) and let \(f: X \to Y\) be the canonical map (TG, I, p. 16). For every i, set \(Y_{i} = f(\mathrm{X}_{i})\) and let \(f_{i}: X_{i} \to Y_{i}\) be the map deduced from f by passing to subspaces. The map \(f_{i}\) is bijective and continuous. If the map \(f_{i}\) is not a homeomorphism, the canonical descent datum on X relative to f is not effective. For an example, cf. TG, I, p. 94, exerc. 15.
\item
  Let B denote the set of points \((x,y)\) of the square \([0,1]^{2}\) satisfying x=0, or x=1, or y=0, or y=1, or there exists \(n\in N\) such that nx=1. Let \(B_{1}\) and \(B_{2}\) the subsets of B consisting of the points \((x,y)\) of B for which \(y\leqslant\frac{1}{2}\) and \(y\geqslant\frac{1}{2}\) respectively.
\end{enumerate}

Construct a B-space \((\mathrm{E}, p)\) which is not a covering but such that the \(B_{1}\) -space \((\mathrm{E}_{\mathrm{B}_{1}}, p_{\mathrm{B}_{1}})\) and the \(B_{2}\) -space \((\mathrm{E}_{\mathrm{B}_{2}}, p_{\mathrm{B}_{2}})\) are coverings.

\begin{enumerate}
\def\labelenumi{\arabic{enumi})}
\setcounter{enumi}{3}
\tightlist
\item
  A topological space is said to be locally simply connected by arcs if every point has a base of neighborhoods that are simply connected by arcs.
\end{enumerate}

\begin{enumerate}
\def\labelenumi{\alph{enumi})}
\item
  Give an example of a topological space that is simply connected by arcs but is not locally simply connected by arcs.
\item
  Let X be a topological space and let G be a discrete group acting properly on X. If X is locally simply connected by arcs, then so is X/G.
\end{enumerate}

\begin{enumerate}
\def\labelenumi{\arabic{enumi})}
\setcounter{enumi}{4}
\item
  Let G be a Lie group acting properly on a loopable and completely regular topological space X. Set Y = X/G and denote \(f: X \to Y\) the canonical map. Prove that the canonical groupoid morphism \(\varpi(X)/G \to \text{Coeg}(f)\) is an isomorphism. (Use LIE, IX, §9, exerc. 17 to adapt the proof of Theorem 3 of IV, p. 403.)
\item
  Let X be a path-connected topological space, let G be a group acting continuously on X, and let x be a point of X such that \(g \cdot x = x\) for all \(g \in G\) . Denote by p the canonical map from X to X/G. Let H be a groupoid and let \(\varphi: \varpi(X) \to H\) a morphism of groupoids. Suppose that for every path \(c\) on \(\mathrm{X}\) and every \(g\in \mathrm{G}\) , we have \(\varphi ([g\cdot c]) = \varphi ([c])\) and that for every loop \(c\in \Omega_x(\mathrm{X})\) , \(\varphi ([c]) = e_{\varphi (x)}\) .
\end{enumerate}

Let \(c_{1}\) and \(c_{2}\) be paths in X such that \(p \circ c_{1}\) and \(p \circ c_{2}\) have the same origin and the same endpoint. Prove that we have \(\varphi([c_{1}]) = \varphi([c_{2}])\) .

\begin{enumerate}
\def\labelenumi{\arabic{enumi})}
\setcounter{enumi}{6}
\tightlist
\item
  Let X be a loopable topological space and let G be a discrete group acting properly on X; we denote by \(m\colon G \times X \to X\) the action of G, \(\mathrm{pr}_2\colon G \times X \to X\) the second projection, and \(f\colon X \to X / G\) the canonical surjection.
\end{enumerate}

\begin{enumerate}
\def\labelenumi{\alph{enumi})}
\item
  Prove that the groupoid morphism \(\varpi(f)\) is surjective.
\item
  Let H be a groupoid and let \(\varphi\colon\varpi(X)\to H\) a morphism of groupoids such that \(\varphi\circ\varpi(m)=\varphi\circ\varpi(\mathrm{pr}_{2})\) .
\end{enumerate}

Let \(c_{1}\) and \(c_{2}\) be paths in X such that the paths \(f \circ c_{1}\) and \(f \circ c_{2}\) in X/G are strictly homotopic. Prove that one has \(\varphi([c_{1}]) = \varphi([c_{2}])\) .

\begin{enumerate}
\def\labelenumi{\alph{enumi})}
\setcounter{enumi}{2}
\tightlist
\item
  Prove that there exists a unique groupoid morphism \(\varphi'\) from \(\varpi(\mathrm{X/G})\) into H such that \(\varphi = \varphi' \circ \varpi(f)\). Deduce from this a new proof of Theorem 3 of IV, p. 403.
\end{enumerate}

\exercisegroup

\begin{enumerate}
\def\labelenumi{\arabic{enumi})}
\tightlist
\item
  Let A be a closed and discrete subset of \(\mathbf{R}^2\) .
\end{enumerate}

\begin{enumerate}
\def\labelenumi{\alph{enumi})}
\item
  For R \textgreater{} 0, let \(D_{R}\) be the set of \(z \in C\) such that \(|z| < R\) ; show that the fundamental group of \(D_{R} - (A \cap D_{R})\) is isomorphic to the free group on \(A \cap D_{R}\) .
\item
  Prove that there exists a homeomorphism \(f\) of \(\mathbf{R}^{2}\) into itself such that \(f(\mathrm{A})\subset\mathbf{N}\times\{0\}\) .
\item
  Show that the fundamental group of \(\mathbf{R}^2 - \mathbf{A}\) is isomorphic to \(\mathrm{F(A)}\) .
\end{enumerate}

\begin{enumerate}
\def\labelenumi{\arabic{enumi})}
\setcounter{enumi}{1}
\item
  Let A be the set of points of \(R^{2}\) with coordinates \((1/n,0)\) , for \(n\in N^{*}\) . Show that the fundamental group of \(R^{2}-A\) has the cardinality of the continuum. In particular, it is not isomorphic to the group F(A).
\item
  Let \(((\mathrm{X}_{i}, x_{i}))_{i \in \mathrm{I}}\) be a family of pointed topological spaces; denote by \((\mathrm{X}, x)\) the wedge of this family (IV, p. 421, example 4). Suppose that, for every \(i \in I\) , the point \(x_{i}\) is closed in \(X_{i}\) and possesses a fundamental system of neighborhoods V such that the group \(\pi_{1}(\mathrm{V}, x_{i})\) is reduced to the identity element. Then, the canonical homomorphism \(*_{i \in \mathrm{I}} \pi_{1}(\mathrm{X}_{i}, x_{i}) \to \pi_{1}(\mathrm{X}, x)\) is an isomorphism.
\item
  Let P be the subspace of \(R^{2}\) defined as in exercise 5, I, p. 146, b). Let a be the point \((0,0,1)\) of \(R^{3}\) and let C be the union of the segments \([a,(x,0)]\) , for \(x \in P\) .
\end{enumerate}

\begin{enumerate}
\def\labelenumi{\alph{enumi})}
\item
  Show that the space C is locally arcwise connected and simply arcwise connected.
\item
  Let D be the union of the space C and its symmetric image -C. Show that the space D is simply connected, but that it is not simply arcwise connected.
\item
  Show that the Poincaré group \(\pi_{1}(D,0)\) has the cardinality of the continuum.
\end{enumerate}

\begin{enumerate}
\def\labelenumi{\arabic{enumi})}
\setcounter{enumi}{4}
\tightlist
\item
  Let P be the subspace of \(R^{2}\) defined as in exercise 5, I, p. 146, b); also denote by \(P^{+}\) (resp. \(P^{-}\) ) the set of points \((x,y)\in\mathrm{P}\) such that \(y\geqslant0\) (resp. \(y\leqslant0\) ). Let a be the point \((0,0,1)\) of \(R^{3}\) and let A be the union of the segment [a,0] and the segments \([a,(2r_{n},0)]\) , for \(n\in N\) . We set \(B^{+}=(P^{+}\times\{0\})\cup A\) , \(B^{-}=(P^{-}\times\{0\})\cup A\) and \(B=(P\times\{0\})\cup A\) .
\end{enumerate}

\begin{enumerate}
\def\labelenumi{\alph{enumi})}
\item
  Prove that the space \(B^{+}\) is loopable and that the Poincaré group \(\pi_{1}(B^{+}, a)\) is a free group on N.
\item
  Show that A is contractible to a.
\item
  Show that the canonical homomorphism from \(\pi_{1}(B^{+},a)*\pi_{1}(B^{-},a)\) into \(\pi_{1}(B,a)\) , deduced from the canonical inclusions of \(B^{+}\) and \(B^{-}\) into B, is not surjective.
\end{enumerate}

\begin{enumerate}
\def\labelenumi{\arabic{enumi})}
\setcounter{enumi}{5}
\tightlist
\item
  Let X be the real line R and let A be a subset of X. Let Y be the space deduced from X by contracting A to a point.
\end{enumerate}

\begin{enumerate}
\def\labelenumi{\alph{enumi})}
\item
  If A is discrete, the Poincaré group of Y is a free group.
\item
  If A has an accumulation point, the Poincaré group of Y has the cardinality of the continuum and is not free.
\item
  If A is the closure of the set of points \(1/n\) , for \(n \in \mathbf{N}^*\) , the space X/A is homeomorphic to the Hawaiian earring.
\end{enumerate}

\begin{enumerate}
\def\labelenumi{\arabic{enumi})}
\setcounter{enumi}{6}
\tightlist
\item
  Let G be a graph, S the set of its vertices and A the set of its oriented edges. We denote by \(|G|\) its geometric realization; we identify S with a subset of \(|G|\) and denote by \(p: I \times A \rightarrow |G|\) the canonical map.
\end{enumerate}

\begin{enumerate}
\def\labelenumi{\alph{enumi})}
\item
  Prove that the space \(|G|\) is compact (resp. locally compact) if and only if the graph G is finite (resp. locally finite).
\item
  Show that the canonical map from S to \(|\mathbf{G}|\) induces a bijection from \(\pi_0(\mathbf{G})\) onto \(\pi_0(|\mathbf{G}|)\) .
\item
  Denote by \(\varpi(|\mathrm{G}|)_\mathrm{S}\) the full subgroupoid of \(\varpi(|\mathrm{G}|)\) with vertex set S. Show that there exists a unique isomorphism of groupoids \(\varpi (\mathrm{G})\to \varpi (|\mathrm{G}|)_{\mathrm{S}}\) which is the identity on the set of vertices and which maps the class of an oriented edge \(a\in \mathrm{A}\) to the class of the path \(t\mapsto p(t,a)\)
\item
  Suppose that the graph G is connected and finite; let \(m = 1 + \text{Card}(S) - \text{Card}(A)\) . Show that the Poincaré group of \(|G|\) is a free group with m generators.
\end{enumerate}

\begin{enumerate}
\def\labelenumi{\arabic{enumi})}
\setcounter{enumi}{7}
\tightlist
\item
  Let X be the quotient space of the space \(R \times \{0,1\}\) by the coarsest equivalence relation identifying \((x,0)\) and \((x,1)\) for all \(x \in R^{*}\) (“real line with doubled origin”).
\end{enumerate}

\begin{enumerate}
\def\labelenumi{\alph{enumi})}
\item
  Prove that \(\pi_1(\mathrm{X},x)\) is isomorphic to \(\mathbf{Z}\) , for every point \(x\in \mathrm{X}\) .
\item
  Let \(\mathrm{X}^{(2)}\) be the quotient space of the space \(X^{2}\) by the action of the symmetric group \(S_{2}\) on the space \(X^{2}\) acting by permutation of coordinates. Show that \(\pi_{1}(\mathrm{X}^{(2)},y)=0\) for all \(y\in\mathrm{X}^{(2)}\) .
\end{enumerate}

\begin{enumerate}
\def\labelenumi{\arabic{enumi})}
\setcounter{enumi}{8}
\tightlist
\item
  Let X be a topological space, let R be an equivalence relation on X and let \(p: X \to X/R\) be the canonical surjection. Let \(x \in X\) .
\end{enumerate}

\begin{enumerate}
\def\labelenumi{\alph{enumi})}
\item
  Suppose that X is connected and locally arcwise connected, that the equivalence classes are connected sets, and that the space X/R is loopable. Show that the homomorphism \(p_{*}:\pi_{1}(X,x)\to\pi_{1}(X/R,p(x))\) is surjective.
\item
  Suppose that X = I and that \(\{0,1\}\) is the only equivalence class which is not reduced to a single element. Show that the homomorphism \(p_{*}\) is not surjective.
\end{enumerate}

\begin{enumerate}
\def\labelenumi{\arabic{enumi})}
\setcounter{enumi}{9}
\tightlist
\item
  Let W be the Warsaw circle, defined in exercise 1, IV, p. 455, whose notations we take up again.
\end{enumerate}

\begin{enumerate}
\def\labelenumi{\alph{enumi})}
\item
  Let R be the equivalence relation on W whose only non-singleton classes are B and S; let \(p: W \rightarrow W/R\) be the canonical surjection. Show that the quotient space W/R is homotopy equivalent, but not homeomorphic, to the circle \(S_{1}\) . Show that the homomorphism \(p_{*}\) is not surjective.
\item
  Let R' be the equivalence relation on W for which B ∪ S is the only equivalence class not reduced to a single element; let \(p'\) : W → W/R' be the canonical surjection. Show that the space W/R' is homeomorphic to \(S_{1}\) , that the homomorphism \(p_{*}\) is not surjective, but that the equivalence classes of R' are connected.
\end{enumerate}

\begin{enumerate}
\def\labelenumi{\arabic{enumi})}
\setcounter{enumi}{10}
\tightlist
\item
  Let X be a path-connected topological space and let \(f: X \to X\) be a homeomorphism. Let T be the quotient space of \(X \times I\) by the finest equivalence relation for which the points \((x,0)\) and \((f(x),1)\) are equivalent, for every \(x \in X\) . We denote \(p: X \times I \to T\) the canonical map; we denote by \(j\colon X\to T\) the map given by \(x\mapsto p(x,1/2)\) . Let a be a point of X, let \(c\in\Lambda_{f(a),a}(X)\) be a path of origin \(f(a)\) and of term \(a\) , and let \(\gamma=[c]\) .
\end{enumerate}

\begin{enumerate}
\def\labelenumi{\alph{enumi})}
\item
  Show that the map from I to T given by \(t \mapsto p(a, \frac{1}{2} - t)\) for \(t \in [0, \frac{1}{2}]\) and \(t \mapsto p(c(2t - 1), \frac{3}{2} - t)\) for \(t \in ]\frac{1}{2}, 1]\) is a loop at \(j(a)\) in T. We denote \(\delta\) its class.
\item
  Let S be the unique group homomorphism from \(\pi_{1}(\mathrm{X},a)*\mathbf{Z}\) into \(\pi_{1}(\mathrm{T},j(a))\) which maps a class \(v\in\pi_{1}(\mathrm{X},a)\) to \(j_{*}(v)\) and the element t=1 of Z to the class \(\delta\). Show that the homomorphism S is surjective and that its kernel is the smallest normal subgroup containing the elements \(vt(\gamma^{-1}f_{*}(v)\gamma)^{-1}t^{-1}\) , for \(v\in\pi_{1}(\mathrm{X},a)\) .
\item
  For \(v \in \pi_{1}(X, a)\) we set \(\varphi(v) = \gamma^{-1} f_{*}(v) \gamma\) ; prove that \(\varphi\) is an automorphism of the group \(\pi_{1}(X, a)\) . Let \(\varphi: \mathbf{Z} \to \operatorname{Aut}(\pi_{1}(X, a))\) be the unique group homomorphism that maps 1 to \(\varphi\) . Construct an isomorphism from the group \(\pi_{1}(T, j(a))\) onto the external semidirect product of Z by \(\pi_{1}(X, a)\) with respect to \(\varphi\) .
\end{enumerate}

\begin{enumerate}
\def\labelenumi{\arabic{enumi})}
\setcounter{enumi}{11}
\tightlist
\item
  \begin{enumerate}
  \def\labelenumii{\alph{enumii})}
  \tightlist
  \item
    Let D be a line in \(R^{3}\) and let \(U = R^{3} - D\) . Prove that U is path-connected and that for every point \(a \in U\) , the group \(\pi_{1}(U, a)\) is isomorphic to Z.
  \end{enumerate}
\end{enumerate}

\begin{enumerate}
\def\labelenumi{\alph{enumi})}
\setcounter{enumi}{1}
\tightlist
\item
  Let P be a plane in \(R^{3}\), let C be a circle contained in P and let \(V = R^{3} - C\) . Prove that V is path-connected and that, for every point \(a \in V\) , the group \(\pi_{1}(V, a)\) is isomorphic to Z. (Apply the van Kampen theorem to the closed covering of \(R^{3}\) formed by the two half-spaces bounded by P; one may also note that the spaces U and V are homeomorphic.)
\end{enumerate}

\begin{enumerate}
\def\labelenumi{\arabic{enumi})}
\setcounter{enumi}{12}
\tightlist
\item
   Let D and D' be two distinct lines of \(R^{3}\), and let \(U = R^{3} - (D \cup D')\) .
\end{enumerate}

\begin{enumerate}
\def\labelenumi{\alph{enumi})}
\item
  Prove that U is path-connected.
\item
  Assume that D and D' are disjoint; prove that, for every point \(a \in U\) , the group \(\pi_{1}(U, a)\) is a free group on two generators.
\item
  Assume that D and D' have a common point; prove that, for every point \(a \in U\) , the group \(\pi_{1}(U, a)\) is a free group on three generators.
\end{enumerate}

\begin{enumerate}
\def\labelenumi{\arabic{enumi})}
\setcounter{enumi}{13}
\tightlist
\item
   Let P be a plane in \(R^{3}\), let C be a circle contained in P and let D be a line in \(R^{3}\) . Let \(U = R^{3} - (C \cup D)\) and let a be a point of U.
\end{enumerate}

\begin{enumerate}
\def\labelenumi{\alph{enumi})}
\item
  Prove that U is path-connected.
\item
  Assume that D does not intersect the convex hull of C. Prove that \(\pi_{1}(U,a)\) is a free group on two generators.
\item
  Assume that D does not meet C but does meet the convex hull of C. Prove that \(\pi_{1}(U,a)\) is isomorphic to \(Z^{2}\) .
\item
  Assume that D intersects C at exactly one point. Prove that \(\pi_{1}(U,a)\) is a free group on two generators.
\item
  Assume that D intersects C at two points. Prove that \(\pi_{1}(U,a)\) is a free group on three generators.
\end{enumerate}

\begin{enumerate}
\def\labelenumi{\arabic{enumi})}
\setcounter{enumi}{14}
\tightlist
\item
  A topological space X is said to have the Phragmén–Brouwer property if it is connected and if, for every pair (A, B) of disjoint closed subsets of X such that \(X-A\) and \(X-B\) are connected, then \(X-(A \cup B)\) is connected.
\end{enumerate}

\begin{enumerate}
\def\labelenumi{\alph{enumi})}
\item
  Let X be a locally connected topological space which has the Phragmén–Brouwer property. Let (A, B) be a pair of disjoint closed subsets of X, and let \((a, b)\) be a pair of points of \(X-(A \cup B)\) which belong to the same connected component of \(X-A\) (resp. of \(X-B\)). Prove that a and b belong to the same connected component of \(X-(A \cup B)\) .
\item
  Let X be a connected and locally connected by arcs topological space that does not have the Phragmén–Brouwer property. Prove that, for every point \(x \in X\), there exists a surjective homomorphism from \(\pi_{1}(X, x)\) into Z.
\item
  Let X be a connected and locally connected by arcs topological space that is simply connected by arcs. Prove that X has the Phragmén-Brouwer property.
\end{enumerate}

d) Let X be a separated and locally connected topological space such that, for every point \(x \in \mathrm{X}\) , the space \(\mathrm{X} - \{x\}\) has the Phragmén-Brouwer property. Let A be a subset of X that is homeomorphic to \([0,1]\) ; prove that \(\mathrm{X} - \mathrm{A}\) is connected. (Let \(c: \mathbf{I} \to \mathrm{A}\) be a homeomorphism. Argue by contradiction and construct a sequence \((\mathrm{J}_n)\) where, for every \(n\) , \(\mathrm{J}_n\) is an interval of \(\mathbf{I}\) of length \(2^{-n}\) , contained in \(\mathrm{J}_{n-1}\) if \(n \geqslant 1\) , such that \(\mathrm{X} - c(\mathrm{J}_n)\) is not connected; let \(x\) be the unique point of X common to all the sets \(c(\mathrm{J}_n)\) Then show that \(\mathrm{X} - \{x\}\) is not connected.)

\begin{enumerate}
\def\labelenumi{\alph{enumi})}
\setcounter{enumi}{4}
\tightlist
\item
  Let \(n\) be an integer \(\geqslant 1\) and let A be a subset of \(\mathbf{S}_n\) which is homeomorphic to \([0,1]\) . Prove that \(\mathbf{S}_n - \mathbf{A}\) is connected.
\end{enumerate}

\begin{enumerate}
\def\labelenumi{\arabic{enumi})}
\setcounter{enumi}{15}
\tightlist
\item
  Let A be a subset of \(S_{2}\) which is homeomorphic to \(S_{1}\) . Let a and b be distinct points of A; we denote by \(A'\) and \(A''\) the connected components of \(A-\{a,b\}\) , and one sets \(X=S_{2}-\{a,b\}\) , \(U=X-A'\) and \(V=X-A''\) .
\end{enumerate}

\begin{enumerate}
\def\labelenumi{\alph{enumi})}
\item
  Show that U and V are connected. (Use Exercise 15.)
\item
  Show that for every point \(x \in U\) the homomorphism from \(\pi_{1}(U, x)\) into \(\pi_{1}(X, x)\) deduced from the inclusion of U in X is trivial.
\item
  Prove that \(S_{2}-A\) has exactly two connected components and that their boundary is equal to A.
\item
  Let S be a subset of \(R^{2}\) which is homeomorphic to \(S_{1}\) ; prove that \(R^{2}-S\) has exactly two connected components and that their boundary is equal to S (Jordan's theorem).
\end{enumerate}

\begin{enumerate}
\def\labelenumi{\arabic{enumi})}
\setcounter{enumi}{16}
\tightlist
\item
  We identify \(S_{1}\) with the set of complex numbers of modulus 1, and the sphere \(S_{3}\) with the subspace of \(C^{2}\) consisting of the points \((z,w)\) such that \(|z|^{2} + |w|^{2} = 1\) . Let B and C be the subspaces of \(S_{3}\) defined by \(|z| \leqslant |w|\) and \(|z| \geqslant |w|\) respectively.
\end{enumerate}

\begin{enumerate}
\def\labelenumi{\alph{enumi})}
\item
  Show that B and C are homeomorphic to \(B_{2} \times S_{1}\) .
\item
  For every pair \((m,n)\) of relatively prime integers, we denote by \(K_{m,n}\) the image of the map from \(S_{1}\) into \(S_{3}\) given by \(u\mapsto(u^{m},u^{n})\) ("torus knot").
\item
  Prove that \(S_{3}-K_{m,n}\) is connected and its fundamental group admits a presentation \(\langle x,y;x^{m}=y^{n}\rangle\) .
\end{enumerate}

d) Prove that there exists a homeomorphism \(\varphi\) of \(\mathbf{S}_3\) into itself such that \(\varphi(\mathrm{K}_{m,n}) = \mathrm{K}_{1,0}\) if and only if \(|m| \leqslant 1\) or \(|n| \leqslant 1\) .

\begin{enumerate}
\def\labelenumi{\arabic{enumi})}
\setcounter{enumi}{17}
\tightlist
\item
  Let \((m,n)\) be a pair of natural integers such that 1 \textless{} m \textless{} n. Let G be the group defined by the presentation \(\langle x,y;x^{m}=y^{n}\rangle\) . Let \(z=x^{m}\) and let C be the subgroup of G generated by z.
\end{enumerate}

\begin{enumerate}
\def\labelenumi{\alph{enumi})}
\item
  Prove that C is contained in the center of G.
\item
  Prove that G/C is isomorphic to the free product \((\mathbf{Z}/m\mathbf{Z})*(\mathbf{Z}/n\mathbf{Z})\) .
\item
  Deduce that C is the center of G.
\item
  Let \((p,q)\) be a pair of natural integers such that 1 \textless{} p \textless{} q. Assume that the group \(G'\) defined by the presentation \(\langle x,y; x^{p} = y^{q}\rangle\) is isomorphic to G. Prove that \((p,q) = (m,n)\) .
\item
  Show that the quotient of G by its derived subgroup is isomorphic to Z.
\end{enumerate}

\begin{enumerate}
\def\labelenumi{\arabic{enumi})}
\setcounter{enumi}{18}
\tightlist
\item
  Let G be a graph; denote by S the set of its vertices and by A the set of its directed edges, and \(|\mathbf{G}|\) its geometric realization. Denote \(j\colon \mathrm{S}\to |\mathrm{G}|\) and \(p\colon \mathbf{I}\times \mathbf{A}\to |\mathbf{G}|\) the canonical maps. We assume that the sets S and A are finite.
\end{enumerate}

 A map \(f\colon|G|\to\mathbf{R}^{n}\) is said to be affine (resp. piecewise affine) if for every edge \(a\in A\) , the map \(t\mapsto f\circ p(a,t)\) is affine (resp. piecewise affine). (Cf. III, p. 331, exerc. 6.)

\begin{enumerate}
\def\labelenumi{\alph{enumi})}
\item
   Prove that there exists an injective and piecewise affine map from \(|G|\) into \(R^{3}\) (Begin by showing that there exists an integer \(n \geqslant 1\) and an injective, piecewise affine map from \(|G|\) into \(R^{n}\) .)
\item
   For every continuous injective map \(f\colon|G|\to R^{n}\), prove that there exists a continuous injective map \(g\colon|G|\to R^{n}\) which is piecewise affine and homotopic to f.
\item
   We assume that the map \((o,t)\colon\mathrm{A}\to\mathrm{S}\times\mathrm{S}\) is injective. Let f be an injective map from S to R. Prove that there exists a unique affine map F from \(|G|\) into \(R^{3}\) which maps every vertex \(s\in S\) onto the point \((f(s),f(s)^{2},f(s)^{3})\) of \(R^{3}\) . Prove that F is injective.
\end{enumerate}

We say that the graph G is planar if there exists a continuous injective map from \(|G|\) into \(R^{2}\) .

d) We assume that the graph \(\mathbf{G}\) is planar and that the map \((o,t)\colon\mathbf{A}\to\mathbf{S}\times\mathbf{S}\) is injective. Let \(f\colon|\mathbf{G}|\to\mathbf{R}^{2}\) be a continuous injective map. Prove that there exists an injective affine map \(g\colon|\mathbf{G}|\to\mathbf{R}^{2}\) which is homotopic to \(f\) .

\begin{enumerate}
\def\labelenumi{\arabic{enumi})}
\setcounter{enumi}{19}
\tightlist
\item
   Let \(f \colon \mathbf{I} \to \mathbf{R}^2\) be a piecewise affine loop whose restriction to [0,1[ is injective; denote \(C = f(\mathbf{I})\) .
\end{enumerate}

\begin{enumerate}
\def\labelenumi{\alph{enumi})}
\item
   Prove that \(\mathbf{R}^2 -\mathbf{C}\) has at most two connected components.
\item
   For \(x \in R^{2} - C\) and \(v \in S_{1}\) , let \(n_{v}(x)\) be the number of connected components of \((\mathbf{R}^{2} - \mathbf{C}) \cap (x + \mathbf{R}_{+}v)\) . Prove that there exists a unique locally constant map \(n: R^{2} - C \to \{0,1\}\) such that \(n(x) \equiv n_{v}(x) \pmod{2}\) for all \(v \in S_{1}\) .
\item
  Deduce that \(R^{2}-C\) has exactly two connected components and that their boundaries are equal to C.
\item
   Let \(g: I \to R^{2}\) be a piecewise affine injective map such that \(\overline{g}^{-1}(C) = \{0,1\}\) ; we set \(P = g(I)\) and \(D = C \cup P\) . Prove that \(C - (C \cap P)\) has two connected components; denote them by \(P_{1}\) and \(P_{2}\) . Prove that \(R^{2} - D\) has three connected components and that their boundaries are equal to C, \(P_{1} \cup P\) and \(P_{2} \cup P\) .
\item
   Let G be a finite planar graph; let S be the set of its vertices and A the set of its oriented edges. Let \(f\colon|G|\to\mathbf{R}^{2}\) be a continuous injective piecewise affine map and let P be its image. Prove that the number of connected components of \(\mathbf{R}^{2}-P\) is equal to \(1+\mathrm{Card(S)}-\mathrm{Card(A)}+\mathrm{Card}(\pi_{0}(\mathrm{G}))\) .
\end{enumerate}

\begin{enumerate}
\def\labelenumi{\arabic{enumi})}
\setcounter{enumi}{20}
\item
  Let K be the graph associated with the quiver whose set of vertices is \(\{1,2,3\} \times \{0,1\}\) , the set of arrows is the set of pairs of the form \(((a,0),(b,1))\) , for \(a,b\in \{1,2,3\}\) and whose source and target maps are respectively derived from the first and second projection (“Kuratowski graph”). Prove that the graph K is not planar.
\item
  Let \(n\) be a natural integer and let \(\mathrm{K}_n\) be a complete graph whose vertex set has cardinality \(n\) (II, p. 220, exerc. 5).
\end{enumerate}

\begin{enumerate}
\def\labelenumi{\alph{enumi})}
\item
  Prove that the graphs \(K_{n}\) , for \(1 \leqslant n \leqslant 4\) , are planar.
\item
  Prove that the graph \(K_{5}\) is not planar.
\end{enumerate}

\begin{enumerate}
\def\labelenumi{\arabic{enumi})}
\setcounter{enumi}{22}
\tightlist
\item
  Let \(u\) and \(v\) be two points of \(\mathbf{S}_1\) such that \(\mathcal{I}(v / u) > 0\) ; one defines a family \((a_j)_{j \in \mathbf{Z} / 4\mathbf{Z}}\) by \(a_1 = u\) , \(a_2 = v\) , \(a_3 = -u\) and \(a_4 = -v\) . Denote by \(\mathrm{Q}_1\) and \(\mathrm{Q}_2\) the two segments \([a_1, a_3]\) and \([a_2, a_4]\) in the unit disk \(\mathbf{B}_2\) and let us set \(\mathrm{Q} = \mathrm{Q}_1 \cup \mathrm{Q}_2\) .
\end{enumerate}

\begin{enumerate}
\def\labelenumi{\alph{enumi})}
\tightlist
\item
  Prove that, for every \(j \in Z/4Z\) , there exists a unique arcwise connected component \(A_{j} \in \pi_{0}(B_{2}-Q)\) whose closure contains \(a_{j}\) and \(a_{j-1}\) . Prove that the union of the family \((A_{j})\) is equal to \(B_{2}-Q\) .
\end{enumerate}

 Let us set \(C = B_{2} \times [-1, 1]\) , \(N = (0, 0, 1)\) and \(S = (0, 0, -1)\) . Let \(f: [-1, 1] \to R_{+}\) be a continuous map such that \(f(-1) = f(1) = 0\) and \(0 < f(t) < 1\) for all \(t \in ]-1; 1[\) . Let us then denote by \(L_{1}\) and \(L_{2}\) the images of the maps from \([-1, 1]\) into C given by \(t \mapsto (ta_{1}, 0)\) and \(t \mapsto (ta_{2}, f(t))\) respectively; finally set \(L = L_{1} \cup L_{2}\) .

\begin{enumerate}
\def\labelenumi{\alph{enumi})}
\setcounter{enumi}{1}
\item
  Prove that C-L is arcwise connected.
\item
  We say that a path \(c\) in \(\mathbf{C} - \mathbf{L}\) with origin \(\mathbf{N}\) and endpoint \(\mathbf{S}\) is straight if there exist juxtaposed paths \(c_1\) and \(c_2\) in \(\mathbf{B}_2\) such that \(c\) is the juxtaposition of the paths given by \(t \mapsto (c_1(t), 1)\) , \(t \mapsto (c_1(1), 1 - 2t)\) and \(t \mapsto (c_2(t), -1)\) . Prove that then \(c_1(1) \not\in \mathbb{Q}\) ; we denote by \(\mathrm{A}_c\) the path-connected component of \(\mathbf{B}_2 - \mathbf{Q}\) which contains this point.
\end{enumerate}

d) Let \(c\) and \(c'\) be paths in \(\mathrm{C} - \mathrm{L}\), with origin \(\mathrm{N}\) and endpoint \(\mathrm{S}\) ; suppose they are straight. Prove that they are strictly homotopic if and only if \(\mathrm{A}_c = \mathrm{A}_{c'}\) .

\begin{enumerate}
\def\labelenumi{\alph{enumi})}
\setcounter{enumi}{4}
\tightlist
\item
  For every \(j \in \mathbf{Z}/4\mathbf{Z}\) , let \(c_j\) be a path in \(C-L\) with origin \(N\) and endpoint \(S\) which is straight and such that \(A_{c_j} = A_j\) . Prove the relation
\end{enumerate}

\[
[ c _ {1} ] [ c _ {4} ] ^ {- 1} = [ c _ {2} ] [ c _ {3} ] ^ {- 1}
\]

in the group \(\pi_1(\mathrm{C} - \mathrm{L},\mathrm{N})\)

\begin{enumerate}
\def\labelenumi{\alph{enumi})}
\setcounter{enumi}{5}
\tightlist
\item
   Prove that the canonical homomorphism from the free group on two generators \(x, y\) into \(\pi_1(\mathrm{C} - \mathrm{L}, \mathrm{N})\) which maps \(x\) and \(y\) to \([c_1][c_2]^{-1}\) and \([c_2][c_3]^{-1}\) respectively is an isomorphism of groups. (Apply the van Kampen theorem to the covering of \(\mathrm{C} - \mathrm{L}\) formed by the eight closed sets delimited by the three vector planes of \(\mathbf{R}^3\) containing two of the points of the set \(\{a_1, a_2, \mathrm{N}\}\) .)
\end{enumerate}

\begin{enumerate}
\def\labelenumi{\arabic{enumi})}
\setcounter{enumi}{23}
\tightlist
\item
   Let m be an integer \(\geqslant 1\) and let \((\theta_{j})_{1\leqslant j\leqslant m}\) be a sequence of real numbers such that \(0\leqslant\theta_{1}<\theta_{2}<\cdots<\theta_{n}<2\pi\) . For \(j\in Z/mZ\) we set \(a_{j}=\exp(i\theta_{k})\) , where k is the unique element of \(\{1,\ldots,n\}\) belonging to class j; we also denote by \(Q_{j}\) the segment \([0,a_{j}]\) of \(B_{2}\) . Finally we set \(Q=Q_{1}\cup\cdots\cup Q_{m}\) .
\end{enumerate}

\begin{enumerate}
\def\labelenumi{\alph{enumi})}
\item
  For \(j \in Z/mZ\) , prove that there exists a unique path-connected component \(A_{j}\) of \(B_{2}-Q\) which contains every point \(b \in S_{1}\) such that \(\mathcal{I}(b/a_{j}) > 0\) and \(\mathcal{I}(b/a_{j+1}) < 0\) . Show that the family \((A_{j})\) is formed by pairwise disjoint sets and that its union is equal to \(B_{2}-Q\) .
\item
   Set \(C = B_{2} \times [-1, 1]\) , \(N = (0, 0, 1)\) , \(S = (0, 0, -1)\) , and \(L = Q \times \{0\}\) . We say that a path in \(C-L\), with origin N and endpoint S, is straight if there exist paths \(c_{1}\) and \(c_{2}\) in \(B_{2}\) , juxtaposed, such that c is the juxtaposition of the paths given by \(t \mapsto (c_{1}(t), 1)\) , \(t \mapsto (c_{1}(1), 1 - 2t)\) and \(t \mapsto (c_{2}(t), -1)\) . Prove that then \(c_{1}(1) \notin Q\) ; we denote by \(A_{c}\) the path-connected component of \(B_{2}-Q\) which contains this point.
\item
   Let c and \(c'\) be paths in \(C-L\), with origin N and endpoint S; suppose they are straight. Prove that they are strictly homotopic if and only if \(A_{c} = A_{c'}\) .
\end{enumerate}

d) For every \(j \in \mathbf{Z}/m\mathbf{Z}\) , let \(c_j\) be a loop in \(C-L\), with origin N and endpoint S, which is straight and such that \(\mathrm{A}_{c_j} = \mathrm{A}_j\) . Prove that the homomorphism from the free group \(\mathrm{F}(\mathbf{Z}/m\mathbf{Z})\) into \(\pi_1(\mathrm{C}-\mathrm{L},\mathrm{N})\) which maps \(x_j\) to \([c_j][c_{j+1}]^{-1}\) , for \(j \in \mathbf{Z}/m\mathbf{Z}\) , is surjective, and that its kernel is the smallest normal subgroup which contains the product \(x_1x_2 \ldots x_m\) .

\begin{enumerate}
\def\labelenumi{\arabic{enumi})}
\setcounter{enumi}{24}
\tightlist
\item
   Set \(C = B_{2} \times [-1, 1]\) . Let K be a closed subspace of \(R^{3}\) and let P be its image under the projection given by \((x, y, z) \mapsto (x, y)\) . We suppose that there exist a finite set I, a family \((p_{i})_{i \in \mathrm{I}}\) of points of \(R^{2}\) , a family \((r_{i})_{i \in \mathrm{I}}\) of strictly positive real numbers and a family \((m_{i})_{i \in \mathrm{I}}\) of natural numbers satisfying the following properties, where \(\varphi_{i}\) denotes the map from \(R^{3}\) into \(R^{3}\) given by \(x \mapsto (p_{i}, 0) + r_{i}x\) and \(C_{i} = \varphi_{i}(C) \) :
\end{enumerate}

- The sets \(C_{i}\) are pairwise disjoint;

 - For every \(i \in \mathrm{I}\) the set \(\mathrm{L}_i = \overline{\varphi}_i^{-1}(\mathrm{C}_i \cap \mathrm{K})\) is the subspace \(\mathrm{L}\) defined in exercise 23 if \(m_i = 0\) and in exercise 24 (where we set \(m = m_i\) ) if \(m_i \geqslant 1\) .

- The set \(\mathrm{K}' = \mathrm{K} - \bigcup_{i} (\mathring{\mathrm{C}}_i \cap \mathrm{K})\) is the union of a finite family of closed segments pairwise disjoint in the plane \(\mathbf{R}^2 \times \{0\}\) .

Let \(\mathrm{N} = (x_{\mathrm{N}}, y_{\mathrm{N}}, z_{\mathrm{N}})\) and \(\mathrm{S} = (x_{\mathrm{S}}, y_{\mathrm{S}}, z_{\mathrm{S}})\) be points of \(R^{3}\) such that \(z_{N} > \sup(r_{i})_{i \in I}\) and \(z_{\mathrm{S}} < -\sup(r_{i})_{i \in I}\) .

\begin{enumerate}
\def\labelenumi{\alph{enumi})}
\tightlist
\item
   We say that a path \(c\) in \(\mathbf{R}^3-K\) with origin N and endpoint S is straight if there exist paths \(c_1\) and \(c_2\) in \(\mathbf{R}^2\) , juxtaposed, such that \(c\) is the juxtaposition of the paths given by \(t \mapsto (c_1(t), z_{\mathrm{N}})\) , \(t \mapsto (c_1(1), (1 - t)z_{\mathrm{N}} + tz_{\mathrm{S}})\) , and \(t \mapsto (c_2(t), z_{\mathrm{S}})\) .
\end{enumerate}

Prove that then the point \(c_1(1)\) does not belong to P; we denote by \(A_c\) the path-connected component of \(\mathbf{R}^2 - \mathbf{P}\) which contains this point. Prove that two straight paths \(c\) and \(c'\) are strictly homotopic if and only if \(A_c = A_{c'}\) .

   For every path-connected component \(\alpha\) of \(\mathbf{R}^2-\mathbf{P}\) , we fix a path \(c(\alpha)\) in \(\mathbf{R}^2-\mathbf{K}\) , with origin N and endpoint S which is straight and such that \(\mathrm{A}_{c(\alpha)} = \alpha\) .

\begin{enumerate}
\def\labelenumi{\alph{enumi})}
\setcounter{enumi}{1}
\item
   Let \(i \in I\) be such that \(m_{i} = 0\) and such that \(\overline{\varphi}_{i}^{-1}(K \cap C_{i})\) is the set L defined in exercise 23, whose notations we take up again; we have \(\overline{\varphi}_{i}^{-1}(P \times \{0\} \cap C_{i}) = Q \times \{0\}\) . For all \(j \in Z/4Z\) , denote by \(\alpha_{i,j}\) the path-connected component of \(R^{2} - P\) which contains the image under \(\varphi_{i}\) of the connected component \(A_{j}\) of \(B_{2} - P\) . Show that we have \([c(\alpha_{i,1})][c(\alpha_{i,4})]^{-1} = [c(\alpha_{i,2})][c(\alpha_{i,3})]^{-1}\) in \(\pi_{1}(R^{3} - K, N)\) .
\item
   Let \(\omega\) be an element of \(\pi_0(\mathbf{R}^2-\mathbf{P})\) . Let \(\lambda_\omega\colon \mathrm{F}(\pi_0(\mathbf{R}^2-\mathbf{P}))\to \pi_1(\mathbf{R}^3-\mathbf{K},\mathbf{N})\) be the unique group homomorphism which, for \(\alpha \in \pi_0(\mathbf{R}^2-\mathbf{P})\) , maps the element \(x_\alpha\) onto \([c(\alpha)][c(\omega)]^{-1}\) .
\end{enumerate}

Prove that the homomorphism \(\lambda_{\omega}\) is surjective and its kernel is the smallest normal subgroup of \(\mathrm{F}(\pi_0(\mathbf{R}^2 - \mathbf{P}))\) which contains the element \(x_{\omega}\) and the elements \(x_{\alpha_{i,1}}^{-1}x_{\alpha_{i,2}}x_{\alpha_{i,3}}^{-1}x_{\alpha_{i,4}}\) , for \(i \in \mathrm{I}\) such that \(m_i = 0\) .

\begin{enumerate}
\def\labelenumi{\alph{enumi})}
\setcounter{enumi}{3}
\item
   For every segment \(s\) belonging to \(\pi_{0}(\mathrm{K}^{\prime})\) , we choose elements \(\alpha_{s}\) and \(\alpha_{s}^{\prime}\) of \(\pi_{0}(\mathbf{R}^{2}-\mathbf{P})\) such that the set of connected components of \(R^{2}-P\) whose closure contains \(s\) is equal to \(\{\alpha_{s},\alpha_{s}^{\prime}\}\) . Let \(\mu\colon\mathrm{F}(\pi_{0}(\mathrm{K}^{\prime}))\to\pi_{1}(\mathbf{R}^{3}-\mathbf{K},\mathbf{N})\) be the unique group homomorphism which, for \(s\in\pi_{0}(\mathrm{K}^{\prime})\) , maps the element \(x_{s}\) onto \([c(\alpha_{s})][c(\alpha_{s}^{\prime})]^{-1}\) . Prove that \(\mu\) is surjective. For \(i\in I\) such that \(m_{i}=0\) and \(j\in\{1,2,3,4\}\) , let \(s_{i,j}\) be the unique segment of \(\pi_{0}(\mathrm{K}^{\prime})\) such that \(\varphi_{i}(a_{j})\) meets \(s_{i,j}\) . Prove that there exist elements \(u_{i},v_{i},w_{i}\) of \(\{-1,1\}\) such that \(\mu(x_{s_{i,4}})=\mu(x_{s_{i,2}}^{w_{i}})\) and \(\mu(x_{s_{i,3}})=\mu(x_{s_{i,2}}^{-u_{i}}x_{s_{i,1}}^{v_{i}}x_{s_{i,2}}^{u_{i}})\) . Show that the kernel of \(\mu\) is the smallest normal subgroup of \(\mathrm{F}(\pi_{0}(\mathrm{K}^{\prime}))\) which contains the elements \(x_{s_{i,3}}x_{s_{i,2}}^{-u_{i}}x_{s_{i,1}}^{-v_{i}}x_{s_{i,2}}^{u_{i}}\) and \(x_{s_{i,4}}x_{s_{i,2}}^{-w_{i}}\) for \(i\in I\) .
\item
  Set
\end{enumerate}

\[
k = \operatorname{Card} (\pi_ {0} (\mathrm{K})) - \operatorname{Card} (\{i \in \mathrm{I} \mid m _ {i} > 0 \}) + \frac {1}{2} \sum_ {i \in \mathrm{I}} m _ {i}.
\]

   Prove that \(k \in N\) . Prove that the quotient of \(\pi_{1}(\mathbf{R}^{3} - \mathbf{K})\) by its derived subgroup is isomorphic to \(Z^{k}\) .

\begin{enumerate}
\def\labelenumi{\arabic{enumi})}
\setcounter{enumi}{25}
\tightlist
\item
  In the numerical plane \(\mathbf{R}^2\) , let \(\mathrm{D}_0\) be the disk with center \((0,0)\) and radius 1 and \(\mathrm{C}_0\) its boundary; we set \(a_0 = (-1,0)\) . Let \(g\) be a natural integer \(\geqslant 1\) , let \(r\) be a strictly positive real number and let \((v_1,\ldots,v_g)\) be a sequence of real numbers such that \(r < \inf (v_{1} + 1,(v_{2} - v_{1}) / 2,\dots,(v_{g} - v_{g - 1}) / 2,1 - v_{g})\) . For all \(j\in \{1,\ldots,g\}\) , we denote by \(\mathrm{D}_j\) the closed disk with center \((0,v_j)\)
\end{enumerate}

and of radius \(r\) and by \(\mathrm{C}_j\) the circle with center \((0,v_j)\) and of radius \(r\) ; we also set \(a_{j}=(-r,v_{j})\) . Let X be the space \(\mathrm{D}=\bigcup_{j=1}^{g}\mathring{\mathrm{D}}_{j}\) .

  For every \(j \in \{1, \ldots, j\}\) , let \(u_j\) be the class in \(\varpi(\mathrm{X})\) of a path with origin \(a_0\) and of term \(a_j\) whose image is the segment \([a_0, a_j]\) ; we also denote by \(c_j \in \pi_1(\mathrm{X}, a_j)\) the class of the loop given by \(t \mapsto (-r \cos(2\pi t), -r \sin(2\pi t) + v_j)\) and \(\gamma_j = u_j c_j u_j^{-1}\) .

\begin{enumerate}
\def\labelenumi{\alph{enumi})}
\item
  Let \(\gamma_{0}\) be the class of the loop in X given by \(t \mapsto (-\cos(2\pi t), -\sin(2\pi t))\) . Show that we have \(\gamma_{0} = \gamma_{1} \ldots \gamma_{g}\) .
\item
  Prove that the unique homomorphism from the free group on g generators \(x_{1},\ldots,x_{j}\) into \(\pi_{1}(X,a_{0})\) which maps \(x_{j}\) to \(\gamma_{j}\) is a group isomorphism. (Apply the van Kampen theorem to the covering of X defined by its intersections with the half-planes \(R_{+}\times R\) and \(R_{-}\times R\) .)
\item
  Let \(f \colon \mathrm{X} \to \mathbf{R}_+\) be a continuous function such that \(f^{-1}(0) = \bigcup_{j=0}^{g} \mathrm{C}_j\) and let Y be the subspace of \(\mathbf{R}^3\) consisting of the points \((x, y, z)\) such that \(z^2 = f(x, y)\) . We denote \(i_+\) (resp. \(i_-\) ) the continuous map from X to Y given by \((x, y) \mapsto (x, y, \sqrt{f(x, y)})\) (resp. \((x, y) \mapsto (x, y, -\sqrt{f(x, y)})\) ). For every path class \(\gamma\) in X, one defines path classes in \(\varpi(\mathrm{Y})\) by \(\gamma^+ = (i_+)_{*}(\gamma)\) and \(\gamma^- = (i_-)_*(\gamma)\) .
\end{enumerate}

For every \(j \in \{1, \ldots, g\}\) , one finally defines an element of \(\pi_1(Y, a_0)\) by \(\delta_j = u_j^+ u_j^-\) . Prove that one has the relations \(\gamma_j^+ \delta_j = \delta_j \gamma_j^-\) for \(1 \leqslant j \leqslant g\) , and \(\gamma_1^+ \ldots \gamma_g^+ = \gamma_1^- \ldots \gamma_g^-\) .

\begin{enumerate}
\def\labelenumi{\alph{enumi})}
\setcounter{enumi}{3}
\tightlist
\item
  Let \(\varphi\) be the unique homomorphism from a free group F with 2g generators \(x_{1},\ldots ,x_{g},y_{1},\ldots ,y_{g}\) into \(\pi_1(\mathrm{Y},a_0)\) which maps \(x_{j}\) to \(\gamma_j^+\) and \(y_{j}\) to \(\delta_{j}\) . Show that the homomorphism \(\varphi\) is surjective and its kernel is the smallest normal subgroup of F that contains the element
\end{enumerate}

\[
(x _ {1} ^ {- 1} y _ {1} x _ {1}) (x _ {2} ^ {- 1} y _ {2} x _ {2}) \dots (x _ {g} ^ {- 1} y _ {g} x _ {g}) (y _ {1} \dots y _ {g}) ^ {- 1}.
\]

(Apply the van Kampen theorem to the covering of Y defined by its intersections with the half-spaces \(R^{2} \times R_{+}\) and \(R^{2} \times R_{-}\) .)

\begin{enumerate}
\def\labelenumi{\alph{enumi})}
\setcounter{enumi}{4}
\tightlist
\item
  Let \(\varphi'\) be the unique homomorphism from a free group F with 2g generators \(x_1, \ldots, x_g, y_1, \ldots, y_g\) into \(\pi_1(Y, a_0)\) which, for every \(j\) , maps \(x_j\) to
\end{enumerate}

\[
(\gamma_ {1} ^ {+} \dots \gamma_ {j - 1} ^ {+}) ^ {- 1} \delta_ {j} (\gamma_ {1} ^ {+} \dots \gamma_ {j - 1} ^ {+})
\]

 and \(y_{j}\) to \(\gamma_{j}^{+}\) . Show that the homomorphism \(\varphi'\) is surjective and its kernel is the smallest normal subgroup of F that contains the element

\[
\left(x _ {1} ^ {- 1} y _ {1} ^ {- 1} x _ {1} y _ {1}\right) \dots \left(x _ {g} ^ {- 1} y _ {g} ^ {- 1} x _ {g} y _ {g}\right).
\]

\begin{enumerate}
\def\labelenumi{\arabic{enumi})}
\setcounter{enumi}{26}
\tightlist
\item
  Let D be the unit disk in C and let X be a subspace of D which is a neighborhood of \(S_{1}\) . Let R be the coarsest equivalence relation in X which, for \(u, v \in S_{1}\) , identifies the points \(f(u)\) and \(f(v)\) if \(f(u) = f(v)\) . Let \(p\colon X \to X / R\) be the canonical surjection. Let \(T\) be a compact topological space and let \(f\colon S_1 \to T\) be a continuous surjective map. We set \(a = f(1)\) and denote by \(\alpha \in \pi_1(T, a)\) the class of the loop \(t \mapsto f(e^{2\pi it})\) in \(T\) .
\end{enumerate}

\begin{enumerate}
\def\labelenumi{\alph{enumi})}
\item
  Prove that there exists a unique map \(\sigma\) from T to X/R such that \(\sigma \circ f = p|_{\mathbf{S}_1}\) . Prove that \(\sigma\) defines a homeomorphism from T onto its image \(p(\mathbf{S}_1)\) .
\item
  Let r be a real number such that 0 \textless{} r \textless{} 1; we suppose that X is the set of \(z \in D\) such that \(r \leqslant |z| \leqslant 1\) . Show that X/R is homeomorphic to the mapping cylinder of f. Deduce that the homomorphism \(\sigma_{*}\) is a group isomorphism from \(\pi_{1}(T, a)\) onto \(\pi_{1}(X/R, p(1))\) .
\item
  Suppose that X = D. Show that X/R is homeomorphic to the mapping cone of f. Deduce that the group homomorphism \(\sigma_{*}\colon\pi_{1}(\mathrm{T},a)\to\pi_{1}(\mathrm{X}/\mathrm{R},p(1))\) is surjective and its kernel is the smallest normal subgroup of \(\pi_{1}(\mathrm{T},a)\) which contains \(\alpha\) .
\item
  Let r be a real number such that 0 \textless{} r \textless{} 1; let \(X_{1}\) be the set of \(z \in X\) such that \(|z| \leqslant r\) and \(X_{2}\) the set of \(z \in D\) such that \(r \leqslant |z| \leqslant 1\) ; we suppose that \(X_{2} \subset X\) . Let \(\beta \in \pi_{1}(X_{2}, r)\) be the class of the loop given by \(t \mapsto re^{2\pi it}\) ; let \(\delta \in \varpi_{p(r), p(1)}\) be the class of the path given by \(t \mapsto p((1 - t)r + t)\) .
\end{enumerate}

  Prove that there exists a unique group homomorphism \(\mu\) from \(\pi_1(\mathrm{T},a)*\pi_1(\mathrm{X}_2,r)\) into \(\pi_1(\mathrm{X} / \mathrm{R},p(1))\) which coincides with the homomorphism \(\sigma_*\) on \(\pi_1(\mathrm{T},a)\) and with the homomorphism \(v\mapsto \delta^{-1}p_{*}(v)\delta\) on \(\pi_1(\mathrm{X}_2,p(r))\) . Prove that \(\mu\) is surjective and its kernel is the smallest normal subgroup containing \(\alpha \beta^{-1}\) .

\begin{enumerate}
\def\labelenumi{\arabic{enumi})}
\setcounter{enumi}{27}
\tightlist
\item
  Let X be a topological space. A homotopy \(\sigma: \mathrm{X} \times \mathbf{I} \to \mathrm{X}\) is said to be an isotopy if, for every \(t \in \mathbf{I}\) , the map \(x \mapsto \sigma(x, t)\) deduced from \(\sigma\) by passage to the subspace is a homeomorphism of X onto itself.
\end{enumerate}

\begin{enumerate}
\def\labelenumi{\alph{enumi})}
\item
  Let m be a natural integer and let \(f: X \to R^{n}\) be a continuous map. Prove that there exists an isotopy whose origin is the identity map of \(X \times R^{n}\) and whose term is the map from \(X \times R^{n}\) into \(X \times R^{n}\) given by \((x, y) \mapsto (x, y - f(x))\) .
\item
  Let A and B be subspaces of X; one says that A is isotopic to B if there exists an isotopy \(\sigma\colon\mathbf{X}\times\mathbf{I}\to\mathbf{X}\) whose origin is the identity map on X and such that \(f(\mathbf{A}\times\{1\})=\mathbf{B}\) .
\end{enumerate}

Prove that the relation “A is isotopic to B” is an equivalence relation on the set of subspaces of X (resp. on the set of closed subspaces of X).

\begin{enumerate}
\def\labelenumi{\alph{enumi})}
\setcounter{enumi}{2}
\tightlist
\item
  Let m and n be natural numbers, and let A be a closed subspace of \(R^{m}\) and let B be a closed subspace of \(R^{n}\) . Assume that A and B are homeomorphic. Prove that the subspaces \(A \times \{0\}\) and \(\{0\} \times B\) of \(R^{m+n}\) are isotopic.
\end{enumerate}

\begin{enumerate}
\def\labelenumi{\arabic{enumi})}
\setcounter{enumi}{28}
\tightlist
\item
  Let A be a closed subset of \(R^{2}\) homeomorphic to a closed subset B of R.
\end{enumerate}

\begin{enumerate}
\def\labelenumi{\alph{enumi})}
\item
  Prove that the subspaces \(A \times \{0\}\) and \(\{(0,0)\} \times B\) of \(R^{3}\) are isotopic.
\item
  We assume that \(B \neq R\) ; prove that \(R^{2}-A\) is path-connected.
\item
  Assume that B = R; prove that \(R^{2}-A\) has exactly two path-connected components and that A is the boundary of each of them.
\item
  Let A be a subset of \(S_{2}\) homeomorphic to \(S_{1}\) ; prove that \(S_{2} - A\) has exactly two path-connected components and that A is the boundary of each of them.
\item
  Let A be a subset of \(R^{2}\) homeomorphic to \(S_{1}\) ; prove that \(R^{2}-A\) has exactly two path-connected components and that A is the boundary of each of them.
\end{enumerate}

\begin{enumerate}
\def\labelenumi{\arabic{enumi})}
\setcounter{enumi}{29}
\tightlist
\item
  Let \(v = (v_{1},\ldots ,v_{n})\colon \mathbf{R}^{n}\to \mathbf{R}^{n}\) be a map of class \(\mathrm{C}^1\) and let \(\delta\) be a real number \(>0\) ; we make the following hypotheses:
\end{enumerate}

- The support of \(v\) is contained in \([- \delta, 1 + \delta] \times \mathring{\mathbf{B}}_{n-1}\) ;

- For every \(x \in [0, 1 - \delta[\) , we have \(v_1(x, 0, \ldots, 0) > 0\) , and \(v_1(1, 0, \ldots, 0) = 0\) ;

- For every \(x \in [0,1]\) , we have \(v_{2}(x,0,\ldots,0) = \cdots = v_{n}(x,0,\ldots,0) = 0\) .

Let \(\Phi\colon\mathbf{R}^{n}\times\mathbf{R}\to\mathbf{R}^{n}\) be the integral flow of the map \((x,t)\mapsto v(x)\) (VAR, 9.1.3); for \(t\in\mathbf{R}\) , let \(\Phi_{t}\) be the map \(x\mapsto\Phi(x,t)\) .

\begin{enumerate}
\def\labelenumi{\alph{enumi})}
\item
  Prove that there exists a real number \(\tau\) such that \(\Phi_{\tau}([0,1] \times \{0\}) \subset [1 - \delta, 1] \times \{0\}\) .
\item
  Let \(j\colon[-\delta,1+\delta]\times\mathbf{B}_{n-1}\to\mathbf{R}^{n}\) be a continuous injective map. For every \(t\in[0,1]\) we set \(\mathrm{A}_{t}=j([t,1]\times\{0\})\) . Prove that, for every \(t\in]0,1[\) , the subspaces \(A_{0}\) and \(A_{t}\) are isotopic.
\item
  Let P be a subset of \(\mathbf{R}^{n}\) homeomorphic to \([0,1]\) which is the union of a finite family of segments. Prove that P is isotopic to a segment.
\end{enumerate}

\begin{enumerate}
\def\labelenumi{\arabic{enumi})}
\setcounter{enumi}{30}
\tightlist
\item
  Let \(f = (f_1, \ldots, f_n) \colon \mathbf{I} \to \mathbf{R}^n\) be a map of class \(\mathbf{C}^1\) such that \(f'(t) \neq 0\) for all \(t \in \mathbf{I}\) .
\end{enumerate}

\begin{enumerate}
\def\labelenumi{\alph{enumi})}
\item
  We assume that \(f_{1}^{\prime}(t) > 0\) for all \(t \in I\) . For all \(j \in \{2, \ldots, n\}\) , prove that there exists a continuous map \(g_{j} \colon R \to R\) such that \(f_{j}(x) = g_{j}(f_{1}(x))\) for all \(x \in I\) . Deduce that \(f(\mathbf{I})\) is isotopic to a segment.
\item
   Prove that there exists a real number \(\delta > 0\) and a continuous injective map \(j\colon [-\delta, 1 + \delta] \times \mathbf{B}_{n-1} \to \mathbf{R}^n\) such that \(j(t, 0) = f(t)\) for all \(t \in \mathbf{I} .\) Deduce from Exercise 30 that, for every \(\tau \in [0, 1[\) the set \(f(\mathbf{I})\) is isotopic to \(f([\tau, 1])\) .
\item
  Prove that \(f(\mathbf{I})\) is isotopic to a segment.
\end{enumerate}

\begin{enumerate}
\def\labelenumi{\arabic{enumi})}
\setcounter{enumi}{31}
\tightlist
\item
  Let A be a closed subset of \(\mathbf{R}^3\) homeomorphic to \([0,1]\) . The endpoints of A are the two points \(p\) of A such that \(\mathrm{A} - \{p\}\) is connected (cf. III, p. 280, th. 2); the other points of A will be called interior.
\end{enumerate}

Let \(p\) be a point of A; one says that A is tame at \(p\) if there exists a neighborhood V of \(p\) in \(\mathbf{R}^3\) and a homeomorphism from V onto the unit ball of \(\mathbf{R}^3\) which maps \(\mathrm{A} \cap \mathrm{V}\) onto a segment.

\begin{enumerate}
\def\labelenumi{\alph{enumi})}
\item
  Let p be a point of A such that A is tame at p. Let \((\mathrm{V}_{n})\) be a decreasing sequence of neighborhoods of p in \(R^{3}\) such that \(\bigcap_{n}\mathrm{V}_{n}=\{p\}\) for all n; for every n, let \(a_{n}\) be a point of \(V_{n}-A\) and denote by \(j_{n}\colon\pi_{1}(\mathrm{V}_{n}-\mathrm{A},a)\to\pi_{1}(\mathrm{V}_{1}-\mathrm{A},a)\) the group homomorphism induced by the inclusion of \(V_{n}\) into \(V_{1}\). If p is an endpoint of A (resp. an interior point), prove that there exists an integer m such that for every integer \(n\geqslant m\) the image of the homomorphism \(j_{n}\) is reduced to the identity element (resp. is abelian).
\item
  Assume that A is tame at every point. Prove that there exists a homeomorphism from \(R^{3}\) onto itself that maps A onto a segment. Deduce that \(R^{3}-A\) is connected and simply connected by arcs.
\end{enumerate}

\begin{enumerate}
\def\labelenumi{\arabic{enumi})}
\setcounter{enumi}{32}
\tightlist
\item
  \begin{enumerate}
  \def\labelenumii{(\arabic{enumii})}
  \setcounter{enumii}{5}
  \tightlist
  \item
    Let L be a subspace of \(\mathbf{R}^3\) as defined in Exercise 23, where one takes \(u = (1 - i)/\sqrt{2}\) and \(v = (1 + i)/\sqrt{2}\) . Let \(a\) and \(r\) be real numbers such that \(0 < 2r < a < 1 - 2r\) . We set \(p_1 = (0, -a)\) , \(p_2 = (0, 0)\) and \(p_3 = (0, a)\) ; \(q_1 = (-1, -a)\) , \(q_2 = (-1, 0)\) , \(q_3 = (-1, a)\) ; \(q_1' = (1, -a)\) , \(q_2' = (1, 0)\) , \(q_3' = (1, a)\) . Let K be the union of the sets \(p_1 + rL\) , \(p_2 + rL\) , \(p_3 + rL\) and the segments \([q_1, p_1 - rv]\) , \([q_2, p_2 - ru]\) , \([q_3, p_3 - ru]\) , \([q_1', p_1 + ru]\) , \([q_2', p_3 + ru]\) , \([q_3', p_3 + rv]\) .
  \end{enumerate}
\end{enumerate}

Let \((x_{n})_{n\in \mathbf{Z}}\) be a strictly increasing family of real numbers such that \(\lim_{n\to -\infty}x_n = -1\) and \(\lim_{n\to +\infty}x_n = 1\) ; for \(n\in \mathbf{Z}\) we set \(z_{n} = y_{n} = 1 - |x_{n}|\) . For \(n\in \mathbf{Z}\) , let \(f_{n}\) be the map from \([-1,1]^{3}\) into \(\mathbf{R}^3\) given by

\[
f _ {n} (x, y, z) = \frac {1 - x}{2} (x _ {n}, y y _ {n}, z z _ {n}) + \frac {1 + x}{2} (x _ {n + 1}, y y _ {n + 1}, z z _ {n + 1})
\]

, and we set \(A_{n}=f_{n}(K)\) and \(B_{n}=f_{n}([-1,1]^{3})\) . Let A be the union of the family \((\mathrm{A}_{n})_{n\in\mathbf{Z}}\) .

\begin{enumerate}
\def\labelenumi{\alph{enumi})}
\item
  Prove that A is homeomorphic to [0, 1].

(6) This example is due to E. ARTIN and R. FOX, “Some wild cells and spheres in three-dimensional space,” Annals of Math. 49 (1948), pp. 979–990.

\item
  For every natural integer \(m\) we set
\end{enumerate}

\[
\mathrm{A} _ {m} ^ {\prime} = \bigcup_ {n <   - m} \mathrm{B} _ {n} \cup \bigcup_ {- m \leqslant n \leqslant m} \mathrm{A} _ {m} \cup \bigcup_ {n > m} \mathrm{B} _ {n}.
\]

Prove that \(R^{3}-A_{m}^{\prime}\) is path-connected and that its fundamental group is generated by elements \(a_{n}, b_{n}, c_{n}\) (for \(n \in Z\) such that \(-m \leqslant n \leqslant m\) ) subject to the relations \(b_{-m} = c_{-m}a_{-m}\) , \(c_{n+1}a_{n+1} = c_{n}c_{n+1}\) , \(c_{n+1}b_{n} = a_{n}c_{n+1}\) , \(b_{n}c_{n+1} = b_{n+1}b_{n}\) (for \(n \in Z\) such that \(-m \leqslant n < m\) ) and \(c_{m}a_{m} = b_{m}\) .

\begin{enumerate}
\def\labelenumi{\alph{enumi})}
\setcounter{enumi}{2}
\tightlist
\item
   Prove that \(R^{3}-A\) is path-connected and that its fundamental group is generated by elements \(a_{n}, b_{n}, c_{n}\) (for \(n \in Z\) ) subject to the relations
\end{enumerate}

\[
b _ {n} = c _ {n} a _ {n}, \quad c _ {n + 1} a _ {n + 1} = c _ {n} c _ {n + 1}, \quad c _ {n + 1} b _ {n} = a _ {n} c _ {n + 1}, \quad b _ {n} c _ {n + 1} = b _ {n + 1} b _ {n}
\]

 for \(n \in Z\) .

d) Prove that there exists a homomorphism from \(\pi_1(\mathbf{R}^3 -\mathbf{A})\) into the group \(\mathfrak{S}_5\) which maps the class of \(c_{n}\) to the permutation \((1,2,3,4,5)\mapsto (2,3,4,5,1)\) if \(n\) is odd and to the permutation \((1,2,3,4,5)\mapsto (4,3,5,2,1)\) if \(n\) is even. Deduce that \(\mathbf{R}^3 -\mathbf{A}\) is not simply connected by arcs.

\begin{enumerate}
\def\labelenumi{\arabic{enumi})}
\setcounter{enumi}{33}
\tightlist
\item
  Let X be a path-connected topological space, let G be a discrete group acting continuously on the left on X. Let M be an open subset of X such that G·M = X. Define the sets S, T, the group F, and the group homomorphism \(\varphi\colon\mathrm{F}\to\mathrm{G}\) as in prop. 10 of I, p. 136. Let \(a\) be a point of M.
\end{enumerate}

\begin{enumerate}
\def\labelenumi{\alph{enumi})}
\item
  Let \(c\colon\mathbf{I}\to\mathbf{X}\) be a loop at \(a\) . Prove that there exists a unique element \(g(c)\) of F satisfying the following property: for every integer \(n\) and every sequence \((s_{1},\ldots,s_{n})\) of elements of S such that \(c([(k-1)/n,k/n])\subset s_{k}\mathrm{M}\) for every integer \(k\in\{1,\ldots,n\}\) , we have \(g(c)=x_{s_{n}^{-1}s_{1}}x_{s_{1}^{-1}s_{2}}\ldots x_{s_{n-1}^{-1}s_{n}}\) .
\item
  Prove that there exists a unique group homomorphism \(\gamma\) of \(\pi_1(\mathrm{X},a)\) into F such that \(\gamma ([c]) = g(c)\) for every loop \(c\) in X at \(a\) .
\item
  Prove that the kernel of \(\varphi\) is equal to the image of the homomorphism \(\gamma\) .
\item
  Prove that the kernel of \(\gamma\) is generated by the union of the images of the groups \(\pi_{1}(\mathrm{M}\cup s\cdot\mathrm{M},a)\) , for \(s\in S\) .
\end{enumerate}

\exercisegroup

\begin{enumerate}
\def\labelenumi{\arabic{enumi})}
\tightlist
\item
   Let K be one of the topological fields R, C, H. Let V be a K-vector space possessing a countable generating system, equipped with its finest topology as a topological vector space (EVT, I, §1, p. 11, cor. 4). For every integer \(n \geqslant 1\) , let \(\mathrm{B}_{n}(\mathrm{V})\) be the subspace of \(V^{n}\) consisting of free sequences and \(\mathrm{G}_{n}(\mathrm{V})\) the set of subspaces of dimension n of V. Denote by \(p: \mathrm{B}_{n}(\mathrm{V}) \to \mathrm{G}_{n}(\mathrm{V})\) the map which associates to a sequence \(v \in \mathrm{B}_{n}(\mathrm{V})\) the vector subspace it generates, and equip the set \(\mathrm{G}_{n}(\mathrm{V})\) with the coarsest topology for which the map p is continuous.
\end{enumerate}

\begin{enumerate}
\def\labelenumi{\alph{enumi})}
\item
   Prove that one defines a right action of the group \(\mathbf{GL}(n,\mathrm{K})\) on \(\mathrm{B}_{n}(\mathrm{V})\) by setting \(v \cdot g = (\sum_{i=1}^{n} v_{i} a_{i,j})_{1 \leqslant j \leqslant n}\) for \(v = (v_{1}, \ldots, v_{n}) \in \mathrm{B}_{n}(\mathrm{V})\) and \(g = (a_{i,j}) \in \mathbf{GL}(n, \mathrm{K})\) .
\item
  Prove that the map p induces a homeomorphism of the quotient space \(\mathrm{B}_{n}(\mathrm{V})/\mathbf{G}\mathbf{L}(n,\mathrm{K})\) onto \(\mathrm{G}_{n}(\mathrm{V})\) . Deduce that the space \(\mathrm{G}_{n}(\mathrm{V})\) is paracompact, and metrizable if V is of finite dimension.
\item
  Prove that \(\mathrm{B}_{n}(\mathrm{V})\) is a principal fiber bundle with group \(\mathbf{GL}(n,\mathrm{K})\) and base \(\mathbf{G}_{n}(\mathrm{V})\) .
\item
  Assume that V is of infinite dimension. Prove that the space \(\mathrm{B}_{n}(\mathrm{V})\) is contractible. Deduce that the space \(\mathrm{G}_{n}(\mathrm{V})\) is a classifying space for the group \(\mathbf{GL}(n,\mathrm{K})\) .
\end{enumerate}

\begin{enumerate}
\def\labelenumi{\arabic{enumi})}
\setcounter{enumi}{1}
\tightlist
\item
  Let H be a separated pre-Hilbert space of infinite dimension. Let \((e_{i})_{i\geqslant0}\) be an orthonormal family of vectors of H which generates a dense subspace of H. For every integer \(n\geqslant1\) , let \(V_{n}(H)\) be the subspace of \(H^{n}\) consisting of sequences \((u_{1},\ldots,u_{n})\) which are orthonormal.
\end{enumerate}

\begin{enumerate}
\def\labelenumi{\alph{enumi})}
\item
  Prove that there exists a unique continuous linear map T: H → H such that \(\mathrm{T}(e_{i}) = e_{i+1}\) for every integer i. Prove that one has \(\|T(x)\| = \|x\|\) for all \(x \in H\) .
\item
  Let H' be the subspace of H formed by the vectors orthogonal to \(e_{i}\) , for i \textless{} n. Prove that the following conditions are equivalent: (i) x and \(\mathrm{T}^{n}(x)\) are linearly dependent; (ii) \(\mathrm{T}(x) = x\) ; (iii) \(x \in H'\) .
\item
  Prove that there exists a unique homotopy \(\sigma\colon\mathrm{V}_{n}(\mathrm{H})\times\mathbf{I}\to\mathrm{V}_{n}(\mathrm{H})\) such that, for every \((u_{1},\ldots,u_{n})\in\mathrm{V}_{n}(\mathrm{H})\) and every \(t\in\mathbf{I}\) , \(\sigma((u_{1},\ldots,u_{n}),t)\) is the sequence deduced from the family \((1-t)(u_{1},\ldots,u_{n})+t(\mathrm{T}^{n}(u_{1}),\ldots,\mathrm{T}^{n}(u_{n}))\) by the orthonormalization process (EVT, V, p. 23, prop. 6).
\item
  Prove that there exists a unique homotopy \(\tau\colon\mathrm{V}_{n}(\mathrm{H}^{\prime})\times\mathbf{I}\to\mathrm{V}_{n}(\mathrm{H})\) such that, for every \((u_{1},\ldots,u_{n})\in\mathrm{V}_{n}(\mathrm{H}^{\prime})\) and every \(t\in\mathbf{I}\) , \(\sigma((u_{1},\ldots,u_{n}),t)\) is the sequence deduced from the family \((1-t)(u_{1},\ldots,u_{n})+t(e_{1},\ldots,e_{n})\) by the orthonormalization process.
\item
  Prove that the space \(V_{n}(H)\) is contractible.
\end{enumerate}

\begin{enumerate}
\def\labelenumi{\arabic{enumi})}
\setcounter{enumi}{2}
\tightlist
\item
  Let H be a separated pre-Hilbert space. Let n be an integer \(\geqslant 1\) ; we denote by \(V_{n}(H)\) the subspace of \(H^{n}\) consisting of orthonormal sequences. Let G be a closed subgroup of \(O(n, \mathbf{R})\) .
\end{enumerate}

\begin{enumerate}
\def\labelenumi{\alph{enumi})}
\tightlist
\item
   Prove that the orthogonal group O(n, R) (LIE, IX, §3, n° 5, p. 22) acts properly and freely on the space \(V_{n}(H)\) by the formula
\end{enumerate}

\[
(u _ {1}, \dots , u _ {n}) \cdot (a _ {i, j}) = \left(\sum_ {i = 1} ^ {n} a _ {i, j} u _ {i}\right).
\]

\begin{enumerate}
\def\labelenumi{\alph{enumi})}
\setcounter{enumi}{1}
\item
  Let G be a closed subgroup of \(O(n,\mathbf{R})\) . Prove that the canonical surjection makes \(V_{n}(\mathrm{H})\) a principal fiber bundle with group G and base \(V_{n}(\mathrm{H})/\mathrm{G}\) (Begin by treating the case where \(G = O(n,\mathbf{R})\) .)
\item
  Prove that the quotient space \(V_{n}(H)/G\) is metrizable. (Endow \(H^{n}\) with its natural structure of pre-Hilbert space and consider the gauge d on \(V_{n}(H)\) given by \(d(u,v)=\inf_{g\in G}\|u-v\cdot g\|\) .)
\item
  Prove that the quotient space \(V_{n}(H)/G\) is a classifying space for G.
\item
  Deduce that every compact Lie group has a classifying space which is a metrizable topological space. (Apply LIE, IX, §9, no. 2, p. 91, th. 1.)
\end{enumerate}

\begin{enumerate}
\def\labelenumi{\arabic{enumi})}
\setcounter{enumi}{3}
\item
  By considering a complex Hilbert space, construct in an analogous manner a metrizable classifying space for the unitary group U(n, C) (LIE, IX, §, no. 4, p. 21)
\item
  \begin{enumerate}
  \def\labelenumii{\alph{enumii})}
  \tightlist
  \item
    Let X be a paracompact topological space. Prove that every principal fiber bundle with base X and group R is trivializable.
  \end{enumerate}
\end{enumerate}

\begin{enumerate}
\def\labelenumi{\alph{enumi})}
\setcounter{enumi}{1}
\tightlist
\item
  Let X be the Alexandroff half-line (TG, IV, p. 49, exerc. 12). Construct a principal fiber bundle with base X and group R which is not trivializable.
\end{enumerate}

\begin{enumerate}
\def\labelenumi{\arabic{enumi})}
\setcounter{enumi}{5}
\tightlist
\item
  Let X be the quotient space of the space \(R \times \{0,1\}\) by the finest equivalence relation for which the points \((t,0)\) and \((t,1)\) are equivalent, for every \(x \in R_{+}^{*}\) . We denote by p: \(R \times \{0,1\} \to X\) the canonical surjection.
\end{enumerate}

\begin{enumerate}
\def\labelenumi{\alph{enumi})}
\item
  Prove that the topological space X is homotopic to a point.
\item
  For which points \(x \in X\) is the pointed space (X, x) contractible?
\item
  Let \(U = p(\mathbf{R} \times \{0\})\) and \(V = p(\mathbf{R} \times \{\mathrm{V}\})\) . Let G be a topological group. Construct a bijection from the set \(\mathcal{C}(\mathbf{R}_{+}^{*}; \mathbf{G})\) of continuous maps from \(R_{+}^{*}\) into G onto the set of isomorphism classes of principal fiber bundles with base X and group G.
\item
  Deduce that there exist principal fiber bundles with base X and group R which are not trivializable.
\end{enumerate}

\begin{enumerate}
\def\labelenumi{\arabic{enumi})}
\setcounter{enumi}{6}
\tightlist
\item
  Let G be a topological group and let E be a principal fiber bundle with base \(S_{2}\) and group G. We denote by \(S_{2}^{+}\) and \(S_{2}^{-}\) the two hemispheres, intersections of \(S_{2}\) and the half-spaces defined by the conditions \(z \geqslant 0\) and \(z \leqslant 0\) respectively; we identify \(S_{2}^{+} \cap S_{2}^{-}\) with \(S_{1}\) .
\end{enumerate}

\begin{enumerate}
\def\labelenumi{\alph{enumi})}
\item
  Prove that the bundle E has a section \(s^{+}\) (resp. \(s^{-}\) ) over \(S_{2}^{+}\) (resp. \(S_{2}^{-}\) ). Prove that there exists a unique continuous map \(f\colon S_{1}\to G\) such that \(s^{+}(x)=s^{-}(x)\cdot f(x)\) for all \(x\in S_{1}\) .
\item
  Prove that the strict homotopy class of the map from \(S_{1}\) into G given by \(x \mapsto c(x)c(1)^{-1}\) does not depend on the choice of sections \(s^{+}\) and \(s^{-}\) ; it is denoted \(c(\mathrm{E})\) .
\item
  Prove that two principal fiber bundles E and \(E'\) , with base \(S_{2}\) and group G, are isomorphic if and only if \(c(\mathrm{E}) = c(\mathrm{E}')\) .
\item
  Prove that, for every class \(c \in \pi_{1}(G, e)\) , there exists a principal fiber bundle E, with base \(S_{2}\) and group G such that \(c(E) = c\) .
\end{enumerate}

\begin{enumerate}
\def\labelenumi{\arabic{enumi})}
\setcounter{enumi}{7}
\tightlist
\item
  Let X be a paracompact topological space and let S(X) be its suspension (I, p. 149, exerc. 1); we denote by \(p\colon X \times [0,1] \to S(X)\) the canonical surjection. Let \(a\) be a point of X.
\end{enumerate}

\begin{enumerate}
\def\labelenumi{\alph{enumi})}
\item
  Let G be a topological group and let E be a principal fiber bundle with group G and base S(X). Prove that the restriction of E to the subspace \(p(\mathrm{X} \times [0,1/2])\) admits a continuous section \(s_{0}\) , and that the restriction of E to the subspace \(p(\mathrm{X} \times [1/2,1])\) admits a continuous section \(s_{1}\) .
\item
  Prove that there exists a unique map \(f\colon \mathrm{X}\to \mathrm{G}\) such that \(s_1(x) = s_0(x)\cdot f(x)\) for all \(x\in \mathrm{X}\) .
\item
  Prove that the strict homotopy class in \([(X,a);(G,e)]\) of the pointed map \(x \mapsto f(x)f(a)^{-1}\) does not depend on the choice of sections \(s_{0}\) and \(s_{1}\) ; it is denoted \(c(\mathrm{E})\) .
\item
  Prove that two principal fiber bundles E and \(E'\) , with base S(X) and group G, are isomorphic if and only if \(c(\mathrm{E}) = c(\mathrm{E}')\) .
\item
  Prove that, for every class \(c \in [(X, a); (G, e)]\) , there exists a principal fiber bundle E, with base S(X) and group G, such that \(c(E) = c\) .
\end{enumerate}

\begin{enumerate}
\def\labelenumi{\arabic{enumi})}
\setcounter{enumi}{8}
\tightlist
\item
  We identify the complex projective line \(\mathbf{P}^{1}(\mathbf{C})\) with the set of lines of \(C^{2}\) . Let E be the subspace of \(\mathbf{C}^{2} \times \mathbf{P}^{1}(\mathbf{C})\) consisting of the pairs \(((u, v), x)\) such that \(v \in x\) and \(|u|^{2} + |v|^{2} = 1\) . We equip E with the action of the group \(S_{1}\) given by \(((u, v), x) \cdot z = ((zu, zv), x)\) .
\end{enumerate}

\begin{enumerate}
\def\labelenumi{\alph{enumi})}
\item
  Prove that the second projection \(\mathrm{pr}_{2}\colon\mathrm{E}\to\mathbf{P}^{1}(\mathbf{C})\) endows E with a structure of principal fiber bundle with base \(\mathbf{P}^{1}(\mathbf{C})\) and group \(S_{1}\) .
\item
  Let \(\varphi\) be a homeomorphism of \(\mathbf{S}_2\) onto \(\mathbf{P}^1 (\mathbf{C})\) ; prove that the group \(\pi_1(\mathbf{S}_1,e)\) is generated by the class \(c(\varphi^{*}\mathrm{E})\) .
\end{enumerate}

\backmatter

\specialchapter{Index of notations}

\(\mathcal{C}_{\mathrm{B}}(\mathrm{X};\mathrm{X}^{\prime})\) (B-morphisms) \ldots.. 2 \(\mathrm{Isom}_{\mathrm{B}}(\mathrm{X};\mathrm{X}^{\prime})\) (B-isomorphisms) \ldots.. 2 \(\mathrm{Aut}_{\mathrm{B}}(\mathrm{X})\) (B-automorphisms) .. 2 \(\mathrm{X}_{b}\) (fiber) \ldots.. 2 \(\mathrm{X}_{\mathrm{A}}\) (induced space) \ldots.. 2 \(\mathrm{X} \times_{\mathrm{B}} \mathrm{X}^{\prime}\) (fiber product) \ldots.. 3 \(\Delta_{\mathrm{X}}\) (diagonal) \ldots.. 4 \(\delta_{\mathrm{X}}\) (diagonal morphism) \ldots.. 4 \(f^{*}(\mathrm{X})\) (base change) .. 5 \(\prod_{\mathrm{B}} \mathrm{X}_{i}\) (fiber product) \ldots.. 5 \(\mathcal{C}_{c}(\mathrm{U};\mathrm{V})\) (continuous maps from U to V) \ldots.. 19 \(\mathcal{S}(\mathrm{A};\mathrm{E})\) (continuous sections) .. 33 \(s \mid \mathrm{U}\) (restriction) \ldots.. 43 \(\underline{\mathcal{F}}(\mathrm{B};\mathrm{X})\) (sheaf of maps) \ldots.. 45 \(\underline{\mathcal{C}}(\mathrm{B};\mathrm{X})\) (sheaf of continuous maps) \ldots.. 45 \(\underline{\mathcal{L}}(\mathrm{B};\mathrm{E}), \underline{\mathcal{L}}(\mathrm{E})\) (sheaf of sections) \ldots.. 45 \(\underline{\mathcal{C}}_{\mathrm{B}}(\mathrm{E};\mathrm{E}^{\prime})\) (sheaf of B-morphisms) \ldots.. 45 \(\underline{\mathcal{Mor}}_{\mathrm{B}}(\mathrm{E};\mathrm{E}^{\prime})\) (sheaf of B-morphisms) \ldots.. 46 \(\underline{\mathcal{Isom}}_{\mathrm{B}}(\mathrm{E};\mathrm{E}^{\prime})\) (sheaf of isomorphisms) \ldots.. 46 \(\underline{\mathcal{C}}^{r}(\mathrm{X};\mathrm{Y})\) (sheaf of maps of class \(C^{r}\) ) \ldots.. 46 \(\underline{\mathfrak{P}}(\mathrm{B})\) (sheaf of subspaces) \ldots.. 46 \(\mathrm{E}_{\mathcal{F}}\) (étale space) \ldots.. 49 \(\widetilde{\mathcal{F}}\) (sheaf associated to a presheaf) 53 \(u_{*}(\mathcal{F})\) (direct image of a (pre)sheaf) 57 \(u^{*}(\mathcal{G})\) (inverse image of a (pre)sheaf) 58 \(\mathcal{G}_{\mathrm{A}}\) (induced sheaf) 59 \(\mathrm{G}^{\circ}\) (opposite group) 91 \(\mathcal{C}_{\mathrm{B}}^{\mathrm{G}}(\mathrm{E};\mathrm{E}^{\prime})\) (morphisms of principal fiber bundles) 92 \(\mathrm{E} \times^{\mathrm{G}} \mathrm{F}\) (associated fiber bundle) . 106 \(\mathrm{Z}_{\mathrm{cont}}^{1}(\mathrm{B},\mathcal{U},\mathrm{G})\) (continuous 1-cocycles) 114 \(\mathrm{H}_{\mathrm{cont}}^{1}(\mathrm{B},\mathcal{U},\mathrm{G})\) (cohomology classes) 116 \(\mathrm{P}(\mathrm{B},\mathrm{G})\) (isomorphism classes of principal bundles) 118 \(\mathrm{P}(\mathrm{B},\mathcal{U},\mathrm{G})\) (isomorphism classes of principal bundles) 118 \(\mathrm{H}_{\mathrm{cont}}^{1}(\mathrm{B},\mathrm{G})\) (cohomology classes) 119 \(\mathcal{C}((\mathrm{X},x);(\mathrm{Y},y))\) (pointed continuous maps) 120 \(p_{\mathrm{A}}\) (gauge of the convex part A) 123 \(\mathrm{S}(\mathrm{X})\) (suspension of the space X) 149\\
Som(C) (vertices of a quiver) 152\\
Fl(C) (arrows of a quiver) . 152

\(o_{C}\) (source map of a quiver) \ldots.. 152 \(t_{C}\) (target map of a quiver) \ldots.. 152 \(\mathrm{Fl}_{a,b}(\mathrm{C}), \mathrm{C}_{a,b}\) (arrows joining a to b) \ldots.. 152 Som(φ) (morphism of quivers) \ldots.. 152 \(\mathrm{Fl}(\varphi)\) (morphism of quivers) \ldots.. 152 \(c*d\) (juxtaposed path) \ldots.. 154 \(\pi_{0}(\mathrm{C})\) (connected components) \ldots.. 155 \(\pi_{0}(\varphi)\) (induced map on connected components) .. 155 \(\overline{f}\) (involution on Fl(C)) \ldots.. 155 Ch(C) (set of paths in a quiver) \ldots.. 160 \(\Omega_{C}\) (category of paths of a quiver) \ldots.. 160 \(f^{-1}\) (inverse of an invertible arrow) \ldots.. 161 \(G_{a}\) (isotropy group) \ldots.. 162 Int(f) (isomorphism of isotropy groups) \ldots.. 162 \(\varphi_{a}\) (homomorphism of isotropy groups) \ldots.. 162 Orb(G) (orbits of a groupoid) \ldots.. 162 Orb(φ) (map deduced from φ by passing to orbits) . 162 \(\mathcal{B}(X,p)\) (permutation groupoid) \ldots.. 163 \(\varinjlim_{i}G_{i}\) (inductive limit of groupoids) \ldots.. 164 X ×s X (groupoid of pairs of a map) \ldots.. 165 \(\varphi^{*}G\) (inverse image groupoid) \ldots.. 166 Ker(φ) (kernel of a groupoid morphism) \ldots.. 166 G/H (quotient groupoid) \ldots{} 170 \(\overline{\varphi}\) (groupoid morphism deduced by passing to the quotient) . 171 \(\varpi_{G}\) (groupoid of path classes) \ldots.. 172 \(\varpi(\varphi)\) (groupoid morphism of path classes) \ldots.. 173 F(A) (free group) \ldots.. 173 Grp(C) (free groupoid) \ldots. 174

G/F (groupoid obtained by contraction of arrows) \ldots{} 176 π₁(G, a) (Poincaré group of a graph at a) \ldots{} 178 u ∼ v (homotopy of groupoid morphisms) \ldots{} 181 Coh(φ, ψ) (cohomotoper) \ldots{} 184 * Gᵢ (free product of groups) i∈I 193 Coeg(φ, ψ) (coequalizer) \ldots{} 199 C(X; Y) (continuous maps from X to Y) \ldots{} 230 {[}X; Y{]} (homotopy classes of continuous maps) . 230 f (homotopy class of a continuous map) \ldots{} 230 C((X, x); (Y, y)) (pointed continuous maps) \ldots{} 232 Sₙ (Euclidean sphere) \ldots{} 234 Cyl(f) (cylinder of a map) \ldots{} 237 X ∪ f B (topological space obtained by attachment) \ldots{} 248 sₓ/A (base point of X/A) \ldots{} 252 Côn(f) (cone of a map) 253 c * d (juxtaposed path) \ldots{} 256 ¯ c (opposite path) \ldots{} 256 Λ(X) (path space) \ldots{} 257 Λₓ(X) (paths with origin x) . 257 λₓ,y(X) (paths with origin x and end y) \ldots{} 257 π₀(X) (path-connected components) \ldots{} 259 π₀(f) (map induced by f on path-connected components) \ldots{} 259 ϖₓ,y(X) (path classes from x to y) \ldots{} 290 Ω(X) (loop space) \ldots{} 290 Ωₓ(X) (loops at x) \ldots{} 290 εₓ (constant loop) \ldots{} 290 ϖ(X) (Poincaré groupoid) 292 π₁(X, x) (fundamental group) 292 ϖ(f) (Poincaré groupoid morphism induced by f) \ldots{} 294 f*(γ) (abbreviation of ϖₓ,y(f)(γ)) \ldots{} 294

\(x \cdot \gamma\) (action of the Poincaré groupoid) \ldots.. 304 \(e_{\mathrm{B}}\) (target map) \ldots.. 342 \(\varepsilon_{\mathrm{B}}\) (target map) \ldots.. 342 \(\pi_{1}(\mathrm{G})\) (Poincaré group of a topological group) \ldots.. 375 \(\mathrm{R}_{\tau} \{\xi z_{1}, z_{2}\}\) (equivalence relation associated with a descent datum) \ldots.. 383

\((Z/R_{\tau},q)\) (quotient space of \((Z,p)\) by a descent datum) \ldots. 384 \(C_{\tau,\tau'}(Z;Z')\) (morphisms compatible with descent data) \ldots. 387 \textbar G\textbar{} (geometric realization of a graph) \ldots. 417 \(\bigvee_{i}(X_{i},x_{i})\) (wedge of pointed topological spaces) \ldots{} 421

\specialchapter{Terminological index}

\indexletter{A}

Adjunction, 59 Admissible action ---, 317 subgroup ---, 315 topology ---, 316 Arcwise continuous map, 267 étale, 28 soft, 64 pointed, 120 relatively connected, 353 separated, 25 strict, 18 topologically étale, 28 universally strict, 20 Tree, 157, 174 maximal, 158 maximal oriented, 178 Edge, oriented edge, 155 Frame, 185, 201, 395, 406 of a covering, 409 Attachment Attachment (topological space obtained by ---), 248

\indexletter{B}

Base of a cone, 253, 426 of a cylinder, 238 Hawaiian earring, 146

Wedge of pointed topological spaces, 421 Target of a path, 154 of an arrow, 151

\indexletter{C}

Quiver, 151 associated with an oriented graph, 156 connected components of a ---, 155 connected, 155 Cayley, 221 Schreier, 221 opposite, 216 product of ---, 153 Square, 7 cartesian, 7 composed, 15 fibred, 4 Category, 160 of paths of a quiver, 160 opposite, 224 Cayley Cayley (graph of ---), 221 Change of base for B-spaces, 4 for principal bundles, 93 Path, 160 in a quiver, 154 in a topological space, 256 juxtaposed, 154, 256

length of a --- in a quiver, 154 opposite, 157, 256 connecting x to y, 256 without backtracking, 157, 173 Paths space of ---, 257 inverse of a --- class, 290 juxtaposable, 154, 256 lifting property of ---, 302 strictly homotopic, 289 Path class in a graph, 172 Cohomology class, 119 Classifying (--- space), 448 Cocycle, 114 cohomology class of a ---, 119 continuous, 114 defined by a family of trivializations, 115 trivial, 118 Cocycles cohomologous, 115 Codimension, 358 Coequalizer, 199 Cohomotoper, 185, 199 Arc-connected component, 259 Composition of juxtaposable path classes, 290 Cone base of a ---, 253, 426 of a topological space, 254 of a map, 253, 425 open of a topological space, 254 vertex of a ---, 254, 425 Connected relatively --- map, 353 topological space arc-connected, 258 topological space locally arc-connected, 260 topological space simply connected, 124 topological space simply arc-connected, 340 Contracted of a topological space, 236 Contractible

topological space --- at a point, 234 locally --- topological space, 346 pointed --- topological space, 234 Contraction, 236 canonical contraction of a cone onto its base, 254 canonical contraction of a cylinder onto its base, 239 topological space obtained by contraction, 252 strong contraction, 237 Convex, 122 Screening of a metric space, 362 Cylinder base of a ---, 238 of a map, 237

\indexletter{D}

Degree of a covering, 73 Loopable (topological space ---), 340 Poincaré half-plane, 150 Diagonal, 4 Descent datum, 382 canonical, 383 effective, 384 morphism compatible with a ---, 387 van Kampen datum, 406, 433

\indexletter{E}

Effective (descent datum ---), 384 Equalizer, 153, 165 Neutral element, 160 Cantor set, 270 of finite-character subsets, 279 pointed, 120 preordered filtered to the right, 164 totally ordered without gaps, 274 Entourage, 272 Complementary structure, 208, 212, 407 complete, 208, 407 basic, 192, 194, 206, 211, 407 Space

classifying, 448 Hopf, 373 of paths, 257 of loops, 290 étalé, 29 Fiber bundle associable to a principal fiber bundle, 107 associated with a principal fiber bundle, 106 locally trivial, 68 principal, 91 principal defined by a cocycle, 115 trivial principal, 92 universal principal, 448 trivial, 68 trivializable, 68 Topological space arcwise connected, 258 Baire, 362 loopable, 340 locally arcwise connected, 260 locally contractible, 346 metrizable, 267 obtained by attaching, 248 obtained by contraction, 252 obtained by identification, 430 pointed, 232 Polish, 364 product, 297 simply connected, 124 simply arcwise connected, 340 Souslin, 362 B-Space, 1 étalé associated with a presheaf, 50 soft étalé, 64 induced, 2 product, 5 quotient, 3 Starred Starred (set --- at a point), 234

\indexletter{F}

Sheaf, 43 associated with a presheaf, 53 constant, 86 of mappings, 45 of continuous mappings, 45

of class C mappings \(^{r}\) , 46 of B-morphisms, 45 of continuous sections, 45 of subspaces, 46 deduced by restriction to an open set, 44 of locally constant maps, 45 of morphisms of principal fiber bundles, 446 of isomorphisms of B-spaces, 46 end, 142 flabby, 144 locally constant, 86 soft, 64 product, 46

Sheaves, isomorphism of ---, 47; morphism of ---, 47.

Locally trivial fibration, 68 principal, 91

Fiber, 2

Fixator, 166

Arrow target of a ---, 151 of a quiver, 151 invertible, 161 opposite, 155 origin of a ---, 151 source of a ---, 151 term of a ---, 151

Composable arrows, 159 family of --- composable, 160 product of a family of --- composable, 160

Functor, 161

Forest, 157, 178 maximal, 157 maximal oriented, 183, 188

\indexletter{G}

Germ, 50 Graph, 155 --- complete, 220 associated with a quiver, 156 combinatorial, 158 Cayley, 221

of Schreier, 221 orientation of a ---, 156 geometric realization of a ---, 417 Directed graph, 156 Directed graph associated with a quiver, 156 Finitely presented group, 359 isotropy group, 162 fundamental group, 292 free group, 173, 193 Poincaré group, 292 of a graph at a vertex, 178 of a topological group, 375 Groupoid, 162 associated with a group, 163 coequalizer, 394 of Poincaré, 292 of path classes of a graph, 172 of pairs, 163, 165 of permutations, 163 fundamental groupoid, 205, 292 inverse image, 166, 183 free groupoid, 174 obtained by contraction of arrows, 176 canonical action of the --- of Poincaré, 304 action of a ---, 167 product, 164, 297 quotient, 170 simply transitive, 162, 174 transitive, 162, 168, 170, 192, 197, 395 Groupoids inductive limit of a family of ---, 164

\indexletter{H}

Homotopic pairs of spaces ---, 235 topological spaces ---, 233 Homotopy, 232 of a pair of spaces, 234 Local homomorphism, 369 continuous homomorphism, 369 Homotopic

morphisms of groupoids ---, 180 Homotopy, 229 canonical morphism associated with the cohomoto- py, 185 class of --- composed, 231 class of --- invertible, 232 class of --- pointed, 232 classes of --- of continuous maps, 230 neutral element up to ---, 373 fixed on a set, 229 inverse up to ---, 233 inverse of a class of ---, 233 juxtaposed, 230 free, 299 origin of a ---, 229 pointed, 231 reciprocal up to ---, 233 relation of ---, 230 pointed relation of ---, 232 strict, 289, 340 term of a ---, 229 Juxtaposable homotopies, 229 extension property of ---, 240 Homotopism, 181 Hopf (--- space), 373

\indexletter{I}

Inverse up to homotopy, 181 of a path class, 290 of an invertible arrow, 161 Isomorphism of B-spaces, 1 of category, 161 of sheaves, 47 of presheaves, 47 local isomorphism, 369 B-Isomorphism, 1

\indexletter{J}

Gauge, 123 Juxtaposable paths ---, 154, 256 homotopies ---, 229 sequence of paths ---, 291 Juxtaposition of paths, 154, 160

\indexletter{L}

Constant loop, 154, 340 in a quiver, 154 in a topological space, 256, 303 at a point x, 290 Loops class of --- at a point x, 290 space of ---, 290 freely homotopic, 299 Width of a train, 264 Composition law in a quiver, 159 associative, 160 Length of a train, 264

\indexletter{M}

Canonical morphism associated with the coequalizer,

canonical morphism associated with the cohomoto- py, 185 compatible with descent data, 387 of B-spaces, 1 of quivers, 152 of categories, 161 of sheaves, 47 of Poincaré groupoids induced by a continuous map,

of groupoids deduced from a mor- phism by passage to the quotient,

of groupoids extending a morphism of quivers, 175 of presheaves, 47 of projection, 48 of principal fiber bundles, 91 diagonal, 4 local of topological groups, 369 B-Morphism, 1 Morphisms deduced from a base equipment,

f-Morphism of principal fiber bundles, 94

N Nielsen-Schreier (theorem of ---), 418 Kernel (of a groupoid morphism), 166 torus knot, 468

\indexletter{O}

object, 161 Admissible operation of the fundamental group,

canonical operation of the fundamental group,

canonical operation of the Poincaré groupoid, 304 of a groupoid, 167, 313, 383 free operation of a group, 213 proper, 287, 349 without monodromy of the Poincaré groupoid, 313 without monodromy of a groupoid, 168 without local monodromy of the Poincaré groupoid, 313 Orbits of a groupoid, 162 of an operation of a groupoid on a set, 168 Orientation, 156 of an edge of a graph, 155 Origin map ---, 257 of a path, 154, 256 of an arrow, 151 of a homotopy, 229

\indexletter{P}

Approachable part, 361 Distinguished part relative to a pair, 197 Meager part, 360 PCV (property ---), 37 Full subquiver, 152 subgroupoid, 293 Base point, 120, 252 Presheaf, 42 deduced by restriction to an open set, 43

direct image, 57 inverse image, 58 separated, 143 Presheaves isomorphism of ---, 47 morphism of ---, 47 Product of sheaves, 46 of principal fiber bundles, 93 bundle, 3 bundle of principal fiber bundles, 93 free product of groups, 193 Projection, 1 of index j, 5 Property (PCV), 37 (VK), 405 of Phragmén-Brouwer, 467 of path lifting, 302, 303, 394 of homotopy extension, 240 Universal property of the inductive limit of groupoids, 164 of cylinders, 238 of free products, 193, 209 of the coequalizer, 199 of the cohomotope, 185, 187 of the product groupoid, 164 of a cone, 254 of a free group, 175 of a product of quivers, 153 Cardinality of the continuum, 366

\indexletter{R}

Refinement of a train, 264 Geometric realization of a graph,

Relator, 372 Equivalence relation associated with a descent datum, 383 Continuous lifting, 2, 33 of paths, 302 property of --- of paths, 302 Restriction of a section of a sheaf, 43 Retract

absolute, 322 absolute neighborhood, 322 of a topological space, 235 Retraction, 235 canonical of a cone with its vertex removed onto its base, 254 canonical of a cylinder onto its base, 238 Covering, 69 associable to a principal covering, 108 finite, 75 Galois, 101, 312 locally finite, 75 principal, 97 product, 88 universal, 120, 343

\indexletter{S}

Schreier Schreier (graph of ---), 221 Continuous section, 2, 33 of a (pre)sheaf, 43 Segment, 122 open, 122 Separated (map ---), 25 Simply transitive (groupoid ---), 162 sum of a family of B-spaces, 2 Vertex of a quiver, 151 of a cone, 254, 255, 425 Source of a path, 154 of an arrow, 151 Subquiver, 152 full, 152 Subcategory, 164 Subsheaf, 44 Subgraph, 155 Admissible subgroup of the fundamental group, 315 Subgroupoid, 165 distinguished, 168 distinguished generated by a set of arrows, 169 generated by a subquiver, 165

generated by an oriented subgraph, 292 full, 193, 293 Subpresheaf, 44 Sphere, 234 Strict map ---, 18 universally strict map ---, 20 Uniform structure, 272 Support of a section of a sheaf,

Suspension of a topological space,

\indexletter{T}

Term map ---, 257 of a path, 154, 256 of an arrow, 151 of a homotopy, 229 Hahn and Mazurkiewicz theorem, 272 Jordan theorem, 468 Nielsen-Schreier theorem, 418 Shelah theorem, 366 van Kampen theorem, 431

Stalk, 50 Admissible topology, 316 Train, 264 width of a ---, 264 length of a ---, 264 refinement of a ---, 264 car of a ---, 264 Transitive (--- groupoid), 162 Transport of structure, 166 Trivializable --- fibered space, 68 --- principal fibered space, 92 Canonical trivialization, 95 of a fibered space, 68 of a principal fibered space, 92

\indexletter{V}

Van Kampen datum of ---, 406, 433 complete equipment deduced from a datum of ---, 407 theorem of ---, 431

W W Car of a train, 264
